% Processed OCR chunk: chunks/05_watson_dna.md
% Scope: Watson/Berry DNA selection only. Compile with a CJK-capable LaTeX engine.
% The parent document should load graphicx.

\providecommand{\ocrimage}[2][]{%
  \IfFileExists{#2}{\includegraphics[#1]{#2}}{%
    \fbox{\parbox{0.78\linewidth}{\centering Missing image reference:\\\texttt{#2}}}}}

\chapter{Beginnings of Genetics: From Mendel to Hitler}
\begin{center}
\large 遗传学的开端：从孟德尔到希特勒
\end{center}

\begin{quote}
\noindent\textbf{English.} My mother, Bonnie Jean, believed in genes. She was proud of her father's Scottish origins, and saw in him the traditional Scottish virtues of honesty, hard work, and thriftiness. She, too, possessed these qualities and felt that they must have been passed down to her from him. His tragic early death meant that her only nongenetic legacy was a set of tiny little girl's kilts he had ordered for her from Glasgow. Perhaps therefore it is not surprising that she valued her father's biological legacy over his material one.

\smallskip
\noindent\textbf{中文。} 我的母亲邦妮·琼相信基因。她为父亲的苏格兰出身感到自豪，也在他身上看到了诚实、勤劳、节俭这些传统的苏格兰美德。她自己也具备这些品质，并觉得它们一定是父亲传给她的。父亲英年早逝，留给她的唯一非遗传遗产，是他从格拉斯哥为她订的一套小女孩穿的小格子裙。也许正因如此，她看重父亲留给她的生物学遗产，胜过物质遗产。
\end{quote}

\begin{quote}
\noindent\textbf{English.} Growing up, I had endless arguments with Mother about the relative roles played by nature and nurture in shaping us. By choosing nurture over nature, I was effectively subscribing to the belief that I could make myself into whatever I wanted to be. I did not want to accept that my genes mattered that much, preferring to attribute my Watson grandmother's extreme fatness to her having overeaten. If her shape was the product of her genes, then I too might have a hefty future. However, even as a teenager, I would not have disputed the evident basics of inheritance, that like begets like. My arguments with my mother concerned complex characteristics like aspects of personality, not the simple attributes that, even as an obstinate adolescent, I could see were passed down over the generations, resulting in "family likeness." My nose is my mother's and now belongs to my son Duncan.

\smallskip
\noindent\textbf{中文。} 在成长过程中，我和母亲无休止地争论：在塑造一个人这件事上，先天与后天各自起多大作用。我站在后天一边，实际上就是认同这样一种信念：我可以把自己塑造成任何想成为的人。我不愿承认自己的基因有那么重要，更愿意把沃森家祖母异常肥胖的体形归因于吃得过多。如果她的身材是基因的产物，那么我自己的未来也可能相当沉重。不过，即便在十几岁时，我也不会否认遗传最显而易见的基本事实：种瓜得瓜，种豆得豆。我和母亲争论的，是性格等复杂特征；至于那些简单属性，即使是固执的少年，我也能看出它们会一代代传下去，造成所谓“家族相似”。我的鼻子来自母亲，如今又长在我儿子邓肯脸上。
\end{quote}

\begin{quote}
\noindent\textbf{English.} Sometimes characteristics come and go within a few generations, but sometimes they persist over many. One of the most famous examples of a long-lived trait is known as the "Hapsburg Lip." This distinctive elongation of the jaw and droopiness to the lower lip---which made the Hapsburg rulers of Europe such a nightmare assignment for generations of court portrait painters---was passed down intact over at least twenty-three generations.

\smallskip
\noindent\textbf{中文。} 有些特征会在短短几代之间出现又消失，有些却能延续许多代。一个最著名的长期存在的性状，就是所谓的“哈布斯堡唇”。这种下颌明显拉长、下唇下垂的特征，使欧洲哈布斯堡统治者成为一代又一代宫廷肖像画家的噩梦；它完整地传递了至少二十三代。
\end{quote}

\begin{figure}[htbp]
\centering
\ocrimage[width=0.72\linewidth]{images/3e4cb4b042e59d2bcb9a019553b80b5387727b74271918f84ee6aa8dac26102c.jpg}
\caption{At age eleven, with my sister Elizabeth and my father, James.\\十一岁时，我和姐姐伊丽莎白以及父亲詹姆斯在一起。}
\end{figure}

\begin{quote}
\noindent\textbf{English.} The Hapsburgs added to their genetic woes by intermarrying. Arranging marriages between different branches of the Hapsburg clan and often among close relatives may have made political sense as a way of building alliances and ensuring dynastic succession, but it was anything but astute in genetic terms. Inbreeding of this kind can result in genetic disease, as the Hapsburgs found out to their cost. Charles II, the last of the Hapsburg monarchs in Spain, not only boasted a prize-worthy example of the family lip---he could not even chew his own food---but was also a complete invalid, and incapable, despite two marriages, of producing children.

\smallskip
\noindent\textbf{中文。} 哈布斯堡家族又通过族内通婚加重了自己的遗传困境。安排哈布斯堡家族不同分支之间，且常常是近亲之间的婚姻，作为建立联盟、确保王朝继承的手段，在政治上也许说得通，但从遗传学上看绝不明智。这样的近亲繁殖可能导致遗传病，哈布斯堡家族为此付出了代价。西班牙哈布斯堡王朝最后一位君主查理二世，不仅拥有堪称“获奖级”的家族唇，甚至连自己的食物都无法咀嚼；他还体弱多病，尽管结过两次婚，却不能生育子女。
\end{quote}

\begin{quote}
\noindent\textbf{English.} Genetic disease has long stalked humanity. In some cases, such as Charles II's, it has had a direct impact on history. Retrospective diagnosis has suggested that George III, the English king whose principal claim to fame is to have lost the American colonies in the Revolutionary War, suffered from an inherited disease, porphyria, which causes periodic bouts of madness. Some historians---mainly British ones---have argued that it was the distraction caused by George's illness that permitted the Americans' against-the-odds military success. While most hereditary diseases have no such geopolitical impact, they nevertheless have brutal and often tragic consequences for the afflicted families, sometimes for many generations. Understanding genetics is not just about understanding why we look like our parents. It is also about coming to grips with some of humankind's oldest enemies: the flaws in our genes that cause genetic disease.

\smallskip
\noindent\textbf{中文。} 遗传病长期纠缠着人类。在某些情况下，例如查理二世，它还直接影响了历史。后来的诊断推测认为，英王乔治三世患有一种遗传病卟啉症，这种病会引发周期性的精神失常；而乔治三世最出名的事，大概就是在美国独立战争中失去了北美殖民地。一些历史学家，主要是英国历史学家，认为正是乔治的疾病造成的分心，使美国人取得了原本胜算不大的军事胜利。多数遗传病并没有这样的地缘政治影响，但它们仍会给患病家庭带来残酷而常常悲惨的后果，有时会延续许多代。理解遗传学，不只是理解我们为什么长得像父母；它还意味着直面人类最古老的敌人之一：导致遗传病的基因缺陷。
\end{quote}

\begin{quote}
\noindent\textbf{English.} Our ancestors must have wondered about the workings of heredity as soon as evolution endowed them with brains capable of formulating the right kind of question. And the readily observable principle that close relatives tend to be similar can carry you a long way if, like our ancestors, your concern with the application of genetics is limited to practical matters like improving domesticated animals, for, say, milk yield in cattle, and plants, for, say, the size of fruit. Generations of careful selection---breeding initially to domesticate appropriate species, and then breeding only from the most productive cows and from the trees with the largest fruit---resulted in animals and plants tailor-made for human purposes. Underlying this enormous unrecorded effort is that simple rule of thumb: that the most productive cows will produce highly productive offspring and from the seeds of trees with large fruit large-fruited trees will grow. Thus, despite the extraordinary advances of the past hundred years or so, the twentieth and twenty-first centuries by no means have a monopoly on genetic insight. Although it was not until 1909 that the British biologist William Bateson gave the science of inheritance a name, genetics, and although the DNA revolution has opened up new and extraordinary vistas of potential progress, in fact the single greatest application of genetics to human well-being was carried out eons ago by anonymous ancient farmers. Almost everything we eat---cereals, fruit, meat, dairy products---is the legacy of that earliest and most far-reaching application of genetic manipulations to human problems.

\smallskip
\noindent\textbf{中文。} 一旦进化赋予我们的祖先足以提出恰当问题的大脑，他们一定就开始思考遗传是如何运作的。近亲往往相似，这个很容易观察到的原则，如果像我们的祖先那样只把遗传学应用于实际事务，譬如改良家畜以提高牛奶产量，或改良植物以增大果实，就已经足以走得很远。一代又一代谨慎的选择，先是繁育适合驯化的物种，随后只让产奶最多的牛和果实最大的树参与繁殖，最终产生了专为人类目的而存在的动植物。在这项庞大而无文字记录的努力背后，是一条朴素经验法则：产量最高的母牛会生出高产后代，大果树的种子会长成结大果的树。因此，尽管过去一百多年进步惊人，二十世纪和二十一世纪绝非独占遗传学洞见。虽然直到1909年，英国生物学家威廉·贝特森才把遗传科学命名为 genetics（遗传学），虽然DNA革命已经打开了前所未有的进步前景，但事实上，对人类福祉影响最大的遗传学应用，是在远古时期由无名的古代农民完成的。我们吃的几乎一切东西，谷物、水果、肉类、乳制品，都是那次最早、也最深远的遗传操控应用于人类问题的遗产。
\end{quote}

\begin{quote}
\noindent\textbf{English.} An understanding of the actual mechanics of genetics proved a tougher nut to crack. Gregor Mendel (1822--1884) published his famous paper on the subject in 1866, and it was ignored by the scientific community for another thirty-four years. Why did it take so long? After all, heredity is a major aspect of the natural world, and, more important, it is readily and universally observable: a dog owner sees how a cross between a brown and black dog turns out, and all parents consciously or subconsciously track the appearance of their own characteristics in their children. One simple reason is that genetic mechanisms turn out to be complicated. Mendel's solution to the problem is not intuitively obvious: children are not, after all, simply a blend of their parents' characteristics. Perhaps most important was the failure by early biologists to distinguish between two fundamentally different processes, heredity and development. Today we understand that a fertilized egg contains the genetic information, contributed by both parents, that determines whether someone will be afflicted with, say, porphyria. That is heredity. The subsequent process, the development of a new individual from that humble starting point of a single cell, the fertilized egg, involves implementing that information. Broken down in terms of academic disciplines, genetics focuses on the information and developmental biology focuses on the use of that information. Lumping heredity and development together into a single phenomenon, early scientists never asked the questions that might have steered them toward the secret of heredity. Nevertheless, the effort had been under way in some form since the dawn of Western history.

\smallskip
\noindent\textbf{中文。} 真正理解遗传学的实际机制，就困难得多。格雷戈尔·孟德尔（1822--1884）在1866年发表了关于这一主题的著名论文，但此后整整三十四年，科学共同体都对它置之不理。为什么要这么久？毕竟遗传是自然界的一个重要方面，更重要的是，它显而易见、人人可见：养狗的人会看到棕狗和黑狗杂交会生出怎样的后代；所有父母也都会有意无意地留意自己的特征在孩子身上出现。一个简单原因是，遗传机制确实复杂。孟德尔给出的解决方案并不符合直觉：孩子终究不是父母特征的简单混合。也许最重要的是，早期生物学家未能区分两个根本不同的过程：遗传与发育。今天我们知道，受精卵含有来自父母双方的遗传信息，这些信息决定了一个人是否会患上卟啉症之类疾病。这就是遗传。随后，从受精卵这一单细胞的谦卑起点发育成一个新个体的过程，则涉及对这些信息的执行。若按学科划分，遗传学关注信息本身，发育生物学关注这些信息如何被使用。早期科学家把遗传和发育混为一个现象，因而从未提出那些本可把他们引向遗传秘密的问题。不过，从西方历史的黎明起，相关努力就已经以某种形式展开了。
\end{quote}

\begin{quote}
\noindent\textbf{English.} The Greeks, including Hippocrates, pondered heredity. They devised a theory of "pangenesis," which claimed that sex involved the transfer of miniaturized body parts: "Hairs, nails, veins, arteries, tendons and their bones, albeit invisible as their particles are so small. While growing, they gradually separate from each other." This idea enjoyed a brief renaissance when Charles Darwin, desperate to support his theory of evolution by natural selection with a viable hypothesis of inheritance, put forward a modified version of pangenesis in the second half of the nineteenth century. In Darwin's scheme, each organ---eyes, kidneys, bones---contributed circulating "gemmules" that accumulated in the sex organs, and were ultimately exchanged in the course of sexual reproduction. Because these gemmules were produced throughout an organism's lifetime, Darwin argued any change that occurred in the individual after birth, like the stretch of a giraffe's neck imparted by craning for the highest foliage, could be passed on to the next generation. Ironically, then, to buttress his theory of natural selection Darwin came to champion aspects of Jean-Baptiste Lamarck's theory of inheritance of acquired characteristics---the very theory that his evolutionary ideas did so much to discredit. Darwin was invoking only Lamarck's theory of inheritance; he continued to believe that natural selection was the driving force behind evolution, but supposed that natural selection operated on the variation produced by pangenesis. Had Darwin known about Mendel's work, although Mendel published his results shortly after \emph{The Origin of Species} appeared and Darwin was never aware of them, he might have been spared the embarrassment of this late-career endorsement of some of Lamarck's ideas.

\smallskip
\noindent\textbf{中文。} 希腊人，包括希波克拉底，也思考过遗传。他们提出了“泛生论”，主张性行为涉及微小化身体部件的转移：“毛发、指甲、静脉、动脉、肌腱以及它们的骨骼，虽然因颗粒极小而不可见；在生长过程中，它们逐渐彼此分离。”十九世纪下半叶，查尔斯·达尔文急切地想为自己的自然选择进化论找到一种可行的遗传假说，于是提出了修订版泛生论，这个观念因此短暂复兴。在达尔文的方案中，每个器官，眼睛、肾脏、骨骼，都会贡献循环于体内的“胚芽粒”；这些胚芽粒聚集到性器官中，最终在有性生殖过程中交换。由于这些胚芽粒在生物个体的一生中持续产生，达尔文认为，出生后个体发生的任何变化，例如长颈鹿为够到最高处树叶而伸长脖子，都可以传给下一代。具有讽刺意味的是，为了支撑自然选择理论，达尔文竟然支持了让-巴蒂斯特·拉马克获得性遗传理论的若干方面，而他的进化思想本来恰恰最有力地削弱了这一理论。达尔文只是借用了拉马克的遗传理论；他仍然相信自然选择是进化的驱动力，只是假设自然选择作用于泛生论所产生的变异。如果达尔文知道孟德尔的工作，虽然孟德尔的成果在《物种起源》发表后不久就出版了，达尔文却从未得知，他也许就不必在晚年因认可拉马克某些观点而尴尬了。
\end{quote}

\begin{quote}
\noindent\textbf{English.} Whereas pangenesis supposed that embryos were assembled from a set of minuscule components, another approach, "preformationism," avoided the assembly step altogether: either the egg or the sperm, exactly which was a contentious issue, contained a complete preformed individual called a homunculus. Development was therefore merely a matter of enlarging this into a fully formed being. In the days of preformationism, what we now recognize as genetic disease was variously interpreted: sometimes as a manifestation of the wrath of God or the mischief of demons and devils; sometimes as evidence of either an excess of or a deficit of the father's "seed"; sometimes as the result of "wicked thoughts" on the part of the mother during pregnancy. On the premise that fetal malformation can result when a pregnant mother's desires are thwarted, leaving her feeling stressed and frustrated, Napoleon passed a law permitting expectant mothers to shoplift. None of these notions, needless to say, did much to advance our understanding of genetic disease.

\smallskip
\noindent\textbf{中文。} 泛生论认为胚胎由一套微小部件组装而成；另一种观点“预成论”则干脆省去了组装步骤：卵子或精子中，究竟是哪一个曾引发争论，已经包含了一个完整预成的微型个体，称为“小人”。发育因此不过是把这个小人放大成完整生命。在预成论盛行的年代，我们今天称为遗传病的现象有各种解释：有时被视为上帝震怒或恶魔作祟的表现；有时被看作父亲“种子”过多或不足的证据；有时则被归因于母亲怀孕期间的“邪念”。拿破仑甚至基于这样一种前提，通过了一项允许孕妇顺手牵羊的法律：如果孕妇的欲望受挫，使她感到压力和沮丧，胎儿可能因此畸形。不用说，这些观念都没有多少助益于我们理解遗传病。
\end{quote}

\begin{quote}
\noindent\textbf{English.} By the early nineteenth century, better microscopes had defeated preformationism. Look as hard as you like, you will never see a tiny homunculus curled up inside a sperm or egg cell. Pangenesis, though an earlier misconception, lasted rather longer---the argument would persist that the gemmules were simply too small to visualize---but was eventually laid to rest by August Weismann, who argued that inheritance depended on the continuity of germ plasm between generations and thus changes to the body over an individual's lifetime could not be transmitted to subsequent generations. His simple experiment involved cutting the tails off several generations of mice. According to Darwin's pangenesis, tailless mice would produce gemmules signifying "no tail," and so their offspring should develop a severely stunted hind appendage or none at all. When Weismann showed that the tail kept appearing after many generations of amputees, pangenesis bit the dust.

\smallskip
\noindent\textbf{中文。} 到十九世纪初，更好的显微镜击败了预成论。无论你多么用力观察，都不可能在精子或卵细胞里看到一个蜷缩着的小人。泛生论虽然是更早的误解，却存活得更久一些，因为支持者可以坚持说胚芽粒实在太小，无法看见。最终奥古斯特·魏斯曼终结了它。他认为，遗传依赖的是代际之间生殖质的连续性，因此个体一生中身体发生的变化不能传给后代。他的简单实验，是连续几代剪掉小鼠的尾巴。按照达尔文的泛生论，无尾小鼠会产生表示“无尾”的胚芽粒，其后代就应当长出极度退化的后肢附属物，或者根本没有尾巴。当魏斯曼证明，许多代被截尾小鼠之后，尾巴仍然照常出现时，泛生论便寿终正寝。
\end{quote}

\begin{figure}[htbp]
\centering
\ocrimage[width=0.58\linewidth]{images/dbdb3499f1e7a9fdf01b306b63c809b52b17d23cc88f3666e7b6956cc8d0f509.jpg}
\caption{Genetics before Mendel: a homunculus, a preformed miniature person imagined to exist in the head of a sperm cell.\\孟德尔之前的遗传学：所谓“小人”，即想象中预先成形、存在于精子头部的微型人。}
\end{figure}

\begin{quote}
\noindent\textbf{English.} Gregor Mendel was the one who got it right. By any standards, however, he was an unlikely candidate for scientific superstardom. Born to a farming family in what is now the Czech Republic, he excelled at the village school and, at twenty-one, entered the Augustinian monastery at Brünn. After proving a disaster as a parish priest---his response to the ministry was a nervous breakdown---he tried his hand at teaching. By all accounts he was a good teacher, but in order to qualify to teach a full range of subjects, he had to take an exam. He failed it. Mendel's father superior, Abbot Napp, then dispatched him to the University of Vienna, where he was to bone up full-time for the retesting. Despite apparently doing well in physics at Vienna, Mendel again failed the exam, and so never rose above the rank of substitute teacher.

\smallskip
\noindent\textbf{中文。} 真正说对的是格雷戈尔·孟德尔。不过无论按什么标准看，他都不像会成为科学巨星的人。孟德尔出生于今天捷克境内的一个农民家庭，在村校成绩出众，二十一岁时进入布尔诺的奥古斯丁修道院。作为堂区神父，他表现得一塌糊涂，面对神职工作的反应是精神崩溃；随后他尝试教书。众人都说他是个好教师，但为了取得教授完整课程的资格，他必须参加考试。他失败了。孟德尔的上级、修道院院长纳普随后派他去维也纳大学，让他全职补习准备重考。尽管他在维也纳的物理学方面似乎学得不错，孟德尔还是再次考试失败，因此从未升到代课教师以上的职位。
\end{quote}

\begin{quote}
\noindent\textbf{English.} Around 1856, at Abbot Napp's suggestion, Mendel undertook some scientific experiments on heredity. He chose to study a number of characteristics of the pea plants he grew in his own patch of the monastery garden. In 1865 he presented his results to the local natural history society in two lectures, and, a year later, published them in the society's journal. The work was a tour de force: the experiments were brilliantly designed and painstakingly executed, and his analysis of the results was insightful and deft. It seems that his training in physics contributed to his breakthrough because, unlike other biologists of that time, he approached the problem quantitatively. Rather than simply noting that crossbreeding of red and white flowers resulted in some red and some white offspring, Mendel actually counted them, realizing that the ratios of red to white progeny might be significant---as indeed they are. Despite sending copies of his article to various prominent scientists, Mendel found himself completely ignored by the scientific community. His attempt to draw attention to his results merely backfired. He wrote to his one contact among the ranking scientists of the day, botanist Karl Nägeli in Munich, asking him to replicate the experiments, and he duly sent off 140 carefully labeled packets of seeds. He should not have bothered. Nägeli believed that the obscure monk should be of service to him, rather than the other way around, so he sent Mendel seeds of his own favorite plant, hawkweed, challenging the monk to re-create his results with a different species. Sad to say, for various reasons, hawkweed is not well suited to breeding experiments such as those Mendel had performed on the peas. The entire exercise was a waste of his time.

\smallskip
\noindent\textbf{中文。} 大约在1856年，在纳普院长的建议下，孟德尔开始进行一些关于遗传的科学实验。他选择研究自己在修道院花园一小块地里种植的豌豆的若干性状。1865年，他在当地自然历史学会用两场讲座报告了结果；一年后，又在学会期刊上发表。那项工作堪称杰作：实验设计精妙、执行细致，结果分析敏锐而娴熟。他在物理学方面的训练似乎促成了突破，因为与当时其他生物学家不同，他以定量方式处理问题。孟德尔并不只是注意到红花和白花杂交会产生一些红花后代和一些白花后代，他真正去数了它们，并意识到红白后代的比例可能重要，事实也正是如此。尽管他把论文副本寄给多位著名科学家，科学共同体却完全无视他。他试图引起注意，反而适得其反。他写信给自己在当时重要科学家中唯一的联系人，慕尼黑植物学家卡尔·内格利，请他重复实验，并寄去140包仔细贴好标签的种子。他本不该费这个劲。内格利认为，应当是这个无名修士为他服务，而不是反过来，于是寄给孟德尔自己最喜欢的植物山柳菊种子，要求这位修士用另一个物种重现结果。遗憾的是，出于多种原因，山柳菊并不适合孟德尔用豌豆所做的那类育种实验。整件事完全浪费了他的时间。
\end{quote}

\begin{quote}
\noindent\textbf{English.} Mendel's low-profile existence as monk-teacher-researcher ended abruptly in 1868 when, on Napp's death, he was elected abbot of the monastery. Although he continued his research---increasingly on bees and the weather---administrative duties were a burden, especially as the monastery became embroiled in a messy dispute over back taxes. Other factors, too, hampered him as a scientist. Portliness eventually curtailed his fieldwork: as he wrote, hill climbing had become "very difficult for me in a world where universal gravitation prevails." His doctors prescribed tobacco to keep his weight in check, and he obliged them by smoking twenty cigars a day, as many as Winston Churchill. It was not his lungs, however, that let him down: in 1884, at the age of sixty-one, Mendel succumbed to a combination of heart and kidney disease.

\smallskip
\noindent\textbf{中文。} 孟德尔作为修士、教师和研究者的低调生活，在1868年戛然而止：纳普去世后，他被选为修道院院长。虽然他继续研究，且越来越多地研究蜜蜂和天气，但行政事务成为沉重负担，尤其是修道院卷入一场混乱的补税争端。其他因素也妨碍了他作为科学家的工作。肥胖最终限制了他的野外活动：正如他所写，在一个“万有引力占统治地位的世界里”，爬山对他已变得“非常困难”。医生开给他烟草来控制体重，他照办了，每天抽二十支雪茄，数量与温斯顿·丘吉尔相当。不过，最终拖垮他的并不是肺。1884年，六十一岁的孟德尔死于心脏病和肾病的共同作用。
\end{quote}

\begin{quote}
\noindent\textbf{English.} Not only were Mendel's results buried in an obscure journal, but they would have been unintelligible to most scientists of the era. He was far ahead of his time with his combination of careful experiment and sophisticated quantitative analysis. Little wonder, perhaps, that it was not until 1900 that the scientific community caught up with him. The rediscovery of Mendel's work, by three plant geneticists interested in similar problems, provoked a revolution in biology. At last the scientific world was ready for the monk's peas.

\smallskip
\noindent\textbf{中文。} 孟德尔的结果不仅埋没在一份不起眼的期刊里，而且对当时大多数科学家来说也难以理解。他把严谨实验与精细的定量分析结合起来，远远超前于时代。因此，直到1900年科学共同体才追上他，也许并不奇怪。三位研究类似问题的植物遗传学家重新发现了孟德尔的工作，引发了一场生物学革命。科学世界终于准备好迎接这位修士的豌豆了。
\end{quote}

\begin{quote}
\noindent\textbf{English.} Mendel realized that there are specific factors---later to be called "genes"---that are passed from parent to offspring. He worked out that these factors come in pairs and that the offspring receives one from each parent.

\smallskip
\noindent\textbf{中文。} 孟德尔认识到，有一些特定因子会从亲代传给子代；这些因子后来被称为“基因”。他弄清楚了这些因子成对存在，而后代从父母双方各获得一个。
\end{quote}

\begin{quote}
\noindent\textbf{English.} Noticing that peas came in two distinct colors, green and yellow, he deduced that there were two versions of the pea-color gene. A pea has to have two copies of the G version if it is to become green, in which case we say that it is GG for the pea-color gene. It must therefore have received a G pea-color gene from both of its parents. However, yellow peas can result both from YY and YG combinations. Having only one copy of the Y version is sufficient to produce yellow peas. Y trumps G. Because in the YG case the Y signal dominates the G signal, we call Y "dominant." The subordinate G version of the pea-color gene is called "recessive." Each parent pea plant has two copies of the pea-color gene, yet it contributes only one copy to each offspring; the other copy is furnished by the other parent. In plants, pollen grains contain sperm cells---the male contribution to the next generation---and each sperm cell contains just one copy of the pea-color gene. A parent pea plant with a YG combination will produce sperm that contain either a Y version or a G one. Mendel discovered that the process is random: 50\% of the sperm produced by that plant will have a Y and 50\% will have a G.

\smallskip
\noindent\textbf{中文。} 孟德尔注意到豌豆有绿色和黄色两种截然不同的颜色，于是推断豌豆颜色基因有两种版本，也就是两种等位基因。一个豌豆若要成为绿色，就必须拥有两个G版本；在这种情况下，我们说它在豌豆颜色基因上是GG。因此，它必然分别从父母双方各获得了一个G版本的豌豆颜色基因。然而，黄色豌豆既可能来自YY组合，也可能来自YG组合。只要有一个Y版本，就足以产生黄色豌豆。Y压过G。因为在YG情况下，Y信号支配G信号，我们称Y为“显性”；而从属的G版本豌豆颜色基因则称为“隐性”。每株亲代豌豆都有两个豌豆颜色基因副本，却只给每个后代贡献一个；另一个副本由另一亲本提供。在植物中，花粉粒含有精细胞，这是给下一代的雄性贡献；每个精细胞只含有豌豆颜色基因的一个副本。带有YG组合的亲代豌豆，会产生含有Y版本或G版本的精子。孟德尔发现这个过程是随机的：该植物产生的精子有50\%带Y，另50\%带G。
\end{quote}

\begin{quote}
\noindent\textbf{English.} Suddenly many of the mysteries of heredity made sense. Characteristics, like the Hapsburg Lip, that are transmitted with a high probability, actually 50\%, from generation to generation are dominant. Other characteristics that appear in family trees much more sporadically, often skipping generations, may be recessive. When a gene is recessive an individual has to have two copies of it for the corresponding trait to be expressed. Those with one copy of the gene are carriers: they do not themselves exhibit the characteristic, but they can pass the gene on. Albinism, in which the body fails to produce pigment so the skin and hair are strikingly white, is an example of a recessive characteristic that is transmitted in this way. Therefore, to be albino you have to have two copies of the gene, one from each parent. This was the case with the Reverend Dr. William Archibald Spooner, who was also---perhaps only by coincidence---prone to a peculiar form of linguistic confusion whereby, for example, "a well-oiled bicycle" might become "a well-boiled icicle." Such reversals would come to be termed "spoonerisms" in his honor. Your parents, meanwhile, may have shown no sign of the gene at all. If, as is often the case, each has only one copy, then they are both carriers. The trait has skipped at least one generation.

\smallskip
\noindent\textbf{中文。} 一下子，遗传的许多谜团都有了道理。像哈布斯堡唇这样以高概率，实际上是50\%，代代相传的性状，是显性性状。其他在家谱中出现得零散得多、常常隔代出现的性状，可能是隐性性状。当一个基因是隐性时，个体必须拥有它的两个副本，相应性状才会表现出来。只拥有一个副本的人是携带者：他们自己不表现这个特征，却可以把该基因传下去。白化病就是这样传递的隐性特征之一；患者身体不能产生色素，因此皮肤和头发呈醒目的白色。所以，要成为白化病患者，你必须拥有该基因的两个副本，一个来自父亲，一个来自母亲。威廉·阿奇博尔德·斯普纳牧师博士就是如此；他还容易出现一种特殊的语言混乱，也许只是巧合，例如把“a well-oiled bicycle”（一辆润滑良好的自行车）说成“a well-boiled icicle”（一根煮得很好的冰柱）。这种音位调换后来便以他的名字称为“斯普纳式口误”。与此同时，你的父母可能完全没有显示出该基因的迹象。如果像常见情况那样，他们各自只有一个副本，那么他们都是携带者。这个性状至少跳过了一代。
\end{quote}

\begin{quote}
\noindent\textbf{English.} Mendel's results implied that things---material objects---were transmitted from generation to generation. But what was the nature of these things?

\smallskip
\noindent\textbf{中文。} 孟德尔的结果暗示，有某些东西，某些物质实体，在一代又一代之间传递。但这些东西的本质是什么？
\end{quote}

\begin{quote}
\noindent\textbf{English.} At about the time of Mendel's death in 1884, scientists using ever-improving optics to study the minute architecture of cells coined the term "chromosome" to describe the long stringy bodies in the cell nucleus. But it was not until 1902 that Mendel and chromosomes came together.

\smallskip
\noindent\textbf{中文。} 大约在孟德尔1884年去世前后，科学家利用日益改进的光学设备研究细胞的微细结构，并创造了“染色体”一词，用来描述细胞核中长而丝状的结构。但直到1902年，孟德尔的因子才与染色体联系到一起。
\end{quote}

\begin{figure}[htbp]
\centering
\ocrimage[width=0.70\linewidth]{images/441c947adb700e4e5337fe0936f94205f10859932393b7b8c9ca50eb6e986016.jpg}
\caption{The human X chromosome, as seen with an electron microscope.\\电子显微镜下的人类X染色体。}
\end{figure}

\begin{quote}
\noindent\textbf{English.} A medical student at Columbia University, Walter Sutton, realized that chromosomes had a lot in common with Mendel's mysterious factors. Studying grasshopper chromosomes, Sutton noticed that most of the time they are doubled up---just like Mendel's paired factors. But Sutton also identified one type of cell in which chromosomes were not paired: the sex cells. Grasshopper sperm have only a single set of chromosomes, not a double set. This was exactly what Mendel had described: his pea plant sperm cells also only carried a single copy of each of his factors. It was clear that Mendel's factors, now called genes, must be on the chromosomes.

\smallskip
\noindent\textbf{中文。} 哥伦比亚大学医学生沃尔特·萨顿意识到，染色体与孟德尔那些神秘因子有许多共同点。萨顿研究蝗虫染色体时注意到，它们大多数时候成对存在，正像孟德尔的成对因子一样。但萨顿还发现了一类染色体不成对的细胞：性细胞。蝗虫精子只有一套染色体，而不是两套。这正是孟德尔描述过的情况：他的豌豆精细胞也只携带每个因子的一个副本。显然，孟德尔的因子，如今称为基因，一定在染色体上。
\end{quote}

\begin{quote}
\noindent\textbf{English.} In Germany Theodor Boveri independently came to the same conclusions as Sutton, and so the biological revolution their work had precipitated came to be called the Sutton--Boveri chromosome theory of inheritance. Suddenly genes were real. They were on chromosomes, and you could actually see chromosomes through the microscope.

\smallskip
\noindent\textbf{中文。} 在德国，西奥多·鲍维里独立得出了与萨顿相同的结论。因此，他们的工作所触发的生物学革命，被称为萨顿--鲍维里染色体遗传理论。突然之间，基因变成了真实存在的东西。它们位于染色体上，而染色体确实可以通过显微镜看到。
\end{quote}

\begin{quote}
\noindent\textbf{English.} Not everyone bought the Sutton--Boveri theory. One skeptic was Thomas Hunt Morgan, also at Columbia. Looking down the microscope at those stringy chromosomes, he could not see how they could account for all the changes that occur from one generation to the next. If all the genes were arranged along chromosomes, and all chromosomes were transmitted intact from one generation to the next, then surely many characteristics would be inherited together. But since empirical evidence showed this not to be the case, the chromosomal theory seemed insufficient to explain the variation observed in nature. Being an astute experimentalist, however, Morgan had an idea how he might resolve such discrepancies. He turned to the fruit fly, \emph{Drosophila melanogaster}, the drab little beast that, ever since Morgan, has been so beloved by geneticists.

\smallskip
\noindent\textbf{中文。} 并非所有人都接受萨顿--鲍维里理论。怀疑者之一是同在哥伦比亚大学的托马斯·亨特·摩尔根。他透过显微镜看着那些丝状染色体，却看不出它们如何能够解释一代到下一代之间发生的所有变化。如果所有基因都沿染色体排列，而所有染色体又完整地从一代传到下一代，那么许多性状理应一起遗传。但经验证据显示事实并非如此，所以染色体理论似乎不足以解释自然界中观察到的变异。不过，摩尔根是一位敏锐的实验学家，他想到也许可以解决这些矛盾。他转向了果蝇，\emph{Drosophila melanogaster}，这只不起眼的小生物自摩尔根以来一直深受遗传学家喜爱。
\end{quote}

\begin{figure}[htbp]
\centering
\ocrimage[width=0.70\linewidth]{images/58fbae006993552b21bb6d1e4a5d19607609d27f70da86419fe2eebeb43e61d7.jpg}
\caption{Notoriously camera-shy T. H. Morgan was photographed surreptitiously while at work in the fly room.\\以怕照相闻名的T. H. 摩尔根，在果蝇室工作时被人偷偷拍了下来。}
\end{figure}

\begin{quote}
\noindent\textbf{English.} In fact, Morgan was not the first to use the fruit fly in breeding experiments---that distinction belonged to a lab at Harvard that first put the critter to work in 1901---but it was Morgan's work that put the fly on the scientific map. \emph{Drosophila} is a good choice for genetic experiments. It is easy to find, as anyone who has left out a bunch of overripe bananas during the summer well knows; it is easy to raise, bananas will do as feed; you can accommodate hundreds of flies in a single milk bottle, and Morgan's students had no difficulty acquiring milk bottles, pinching them at dawn from doorsteps in their Manhattan neighborhood; and it breeds and breeds and breeds, a whole generation taking about ten days, with each female laying several hundred eggs. Starting in 1907 in a famously squalid, cockroach-infested, banana-stinking lab that came to be known affectionately as the "fly room," Morgan and his students, "Morgan's boys" as they were called, set to work on fruit flies.

\smallskip
\noindent\textbf{中文。} 事实上，摩尔根并不是第一个把果蝇用于育种实验的人；这一荣誉属于哈佛的一个实验室，那里在1901年率先让这种小虫投入工作。但正是摩尔根的研究让果蝇登上了科学地图。果蝇是遗传实验的好材料。它容易找到，夏天把一串过熟香蕉放在外面的人都知道这一点；它容易饲养，香蕉就能当饲料；一个牛奶瓶里可以容纳数百只果蝇，而摩尔根的学生弄到牛奶瓶毫不费力，他们清晨从曼哈顿邻里门前顺走即可；它还不停繁殖，一代约十天，每只雌蝇产卵数百枚。1907年起，在一个出了名肮脏、蟑螂横行、香蕉味刺鼻，后来被亲切称为“果蝇室”的实验室里，摩尔根和他的学生们，人称“摩尔根的男孩们”，开始研究果蝇。
\end{quote}

\begin{quote}
\noindent\textbf{English.} Unlike Mendel, who could rely on the variant strains isolated over the years by farmers and gardeners---yellow peas as opposed to green ones, wrinkled skin as opposed to smooth---Morgan had no menu of established genetic differences in the fruit fly to draw upon. And you cannot do genetics until you have isolated some distinct characteristics to track through the generations. Morgan's first goal therefore was to find "mutants," the fruit fly equivalents of yellow or wrinkled peas. He was looking for genetic novelties, random variations that somehow simply appeared in the population.

\smallskip
\noindent\textbf{中文。} 孟德尔可以依靠农民和园丁多年分离出来的变异品系，例如黄豌豆相对于绿豌豆、皱皮相对于光滑；摩尔根却没有一份果蝇既有遗传差异的菜单可用。而在分离出某些明确特征并能跨代追踪之前，你无法做遗传学。因此，摩尔根的首要目标是寻找“突变体”，也就是果蝇中相当于黄豌豆或皱豌豆的东西。他寻找的是遗传新奇性，是种群中不知怎么就随机出现的变异。
\end{quote}

\begin{quote}
\noindent\textbf{English.} One of the first mutants Morgan observed turned out to be one of the most instructive. While normal fruit flies have red eyes, these had white ones. And he noticed that the white-eyed flies were typically male. It was known that the sex of a fruit fly---or, for that matter, the sex of a human---is determined chromosomally: females have two copies of the X chromosome, whereas males have one copy of the X and one copy of the much smaller Y. In light of this information, the white-eye result suddenly made sense: the eye-color gene is located on the X chromosome and the white-eye mutation, W, is recessive. Because males have only a single X chromosome, even recessive genes, in the absence of a dominant counterpart to suppress them, are automatically expressed. White-eyed females were relatively rare because they typically had only one copy of W, so they expressed the dominant red eye color. By correlating a gene---the one for eye color---with a chromosome, the X, Morgan, despite his initial reservations, had effectively proved the Sutton--Boveri theory. He had also found an example of "sex-linkage," in which a particular characteristic is disproportionately represented in one sex.

\smallskip
\noindent\textbf{中文。} 摩尔根最早观察到的突变体之一，后来证明极富启发性。正常果蝇有红眼，而这些果蝇是白眼。他还注意到，白眼果蝇通常是雄性。人们已经知道，果蝇的性别，事实上人类的性别也是如此，由染色体决定：雌性有两条X染色体，雄性有一条X染色体和一条小得多的Y染色体。有了这一信息，白眼结果立刻变得有道理：眼色基因位于X染色体上，白眼突变W是隐性的。由于雄性只有一条X染色体，即便是隐性基因，只要没有显性对应基因压制，也会自动表现出来。白眼雌蝇相对罕见，因为它们通常只有一个W副本，因此表现显性红眼。摩尔根把一个基因，也就是眼色基因，与一条染色体X联系起来，尽管他最初有所保留，却实际上证明了萨顿--鲍维里理论。他还发现了“伴性遗传”的例子，即某一特定性状在某一性别中不成比例地出现。
\end{quote}

\begin{quote}
\noindent\textbf{English.} Like Morgan's fruit flies, Queen Victoria provides a famous example of sex-linkage. On one of her X chromosomes, she had a mutated gene for hemophilia, the "bleeding disease" in whose victims proper blood clotting fails to occur. Because her other copy was normal, and the hemophilia gene is recessive, she herself did not have the disease. But she was a carrier. Her daughters did not have the disease either; evidently each possessed at least one copy of the normal version. But Victoria's sons were not all so lucky. Like all males, fruit fly males included, each had only one X chromosome; this was necessarily derived from Victoria, since a Y chromosome could have come only from Prince Albert, Victoria's husband. Because Victoria had one mutated copy and one normal copy, each of her sons had a 50--50 chance of having the disease. Prince Leopold drew the short straw: he developed hemophilia, and died at thirty-one, bleeding to death after a minor fall. Two of Victoria's daughters, Princesses Alice and Beatrice, were carriers, having inherited the mutated gene from their mother. They each produced carrier daughters and sons with hemophilia. Alice's grandson Alexis, heir to the Russian throne, had hemophilia, and would doubtless have died young had the Bolsheviks not gotten to him first.

\smallskip
\noindent\textbf{中文。} 与摩尔根的果蝇一样，维多利亚女王也提供了一个著名的伴性遗传例子。她的一条X染色体上带有血友病的突变基因；血友病又称“出血病”，患者无法正常凝血。由于她的另一个副本正常，而血友病基因是隐性的，她本人并未患病。但她是携带者。她的女儿们也没有患病，显然每个人至少有一个正常版本。可维多利亚的儿子们并非都如此幸运。像所有雄性一样，包括雄性果蝇在内，每个男性只有一条X染色体；这条X必然来自维多利亚，因为Y染色体只能来自维多利亚的丈夫阿尔伯特亲王。由于维多利亚有一个突变副本和一个正常副本，她每个儿子患病的机会都是一半对一半。利奥波德王子抽中了下签：他患上血友病，三十一岁时因一次小摔倒后流血不止而死。维多利亚的两个女儿，爱丽丝公主和比阿特丽斯公主，也是携带者，她们从母亲那里继承了突变基因。她们各自生下携带者女儿和血友病儿子。爱丽丝的孙子、俄国王位继承人阿列克谢患有血友病；若不是布尔什维克先下了手，他无疑也会早夭。
\end{quote}

\begin{quote}
\noindent\textbf{English.} Morgan's fruit flies had other secrets to reveal. In the course of studying genes located on the same chromosome, Morgan and his students found that chromosomes actually break apart and re-form during the production of sperm and egg cells. This meant that Morgan's original objections to the Sutton--Boveri theory were unwarranted: the breaking and re-forming---"recombination," in modern genetic parlance---shuffles gene copies between members of a chromosome pair. This means that, say, the copy of chromosome 12 I got from my mother, the other of course comes from my father, is in fact a mix of my mother's two copies of chromosome 12, one of which came from her mother and one from her father. Her two 12s recombined---exchanged material---during the production of the egg cell that eventually turned into me. Thus my maternally derived chromosome 12 can be viewed as a mosaic of my grandparents' 12s. Of course, my mother's maternally derived 12 was itself a mosaic of her grandparents' 12s, and so on.

\smallskip
\noindent\textbf{中文。} 摩尔根的果蝇还有其他秘密可揭示。在研究位于同一染色体上的基因时，摩尔根和学生发现，染色体在精子和卵细胞形成过程中实际上会断裂并重新连接。这意味着摩尔根最初对萨顿--鲍维里理论的反对并不成立：断裂和重组，用现代遗传学术语说就是“重组”，会在一对染色体成员之间洗牌式地交换基因副本。也就是说，比如我从母亲那里得到的第12号染色体副本，另一条当然来自父亲，事实上是母亲两条12号染色体的混合体；其中一条来自她的母亲，另一条来自她的父亲。在形成最终变成我的那枚卵细胞时，她的两条12号染色体发生了重组，也就是交换了物质。因此，我母源性的12号染色体可以看作外祖父母两条12号染色体的镶嵌体。当然，我母亲从她母亲那里得到的12号染色体本身，又是她的祖父母两条12号染色体的镶嵌体，如此上溯。
\end{quote}

\begin{quote}
\noindent\textbf{English.} Recombination permitted Morgan and his students to map out the positions of particular genes along a given chromosome. Recombination involves breaking and re-forming chromosomes. Because genes are arranged like beads along a chromosome string, a break is statistically much more likely to occur between two genes that are far apart, with more potential break points intervening, on the chromosome than between two genes that are close together. If, therefore, we see a lot of reshuffling for any two genes on a single chromosome, we can conclude that they are a long way apart; the rarer the reshuffling, the closer the genes likely are. This basic and immensely powerful principle underlies all of genetic mapping. One of the primary tools of scientists involved in the Human Genome Project and of researchers at the forefront of the battle against genetic disease was thus developed all those years ago in the filthy, cluttered Columbia fly room. Each new headline in the science section of the newspaper these days along the lines of "Gene for Something Located" is a tribute to the pioneering work of Morgan and his boys.

\smallskip
\noindent\textbf{中文。} 重组使摩尔根和学生能够绘制特定基因在给定染色体上的位置图。重组涉及染色体的断裂与重新连接。由于基因像珠子一样沿染色体这根线排列，两个相距很远的基因之间有更多潜在断点，因此断裂在它们之间发生的统计概率，远高于在两个相邻基因之间发生的概率。所以，如果我们看到同一条染色体上任意两个基因发生大量重新组合，就可以断定它们相距很远；重组越少，基因就越可能靠得近。这个基本而极其有力的原则，是全部遗传作图的基础。后来人类基因组计划科学家以及抗击遗传病前沿研究者的一项主要工具，竟是在多年前哥伦比亚大学那间肮脏杂乱的果蝇室里发展出来的。如今报纸科学版每出现一条类似“某某基因定位成功”的新标题，都是对摩尔根和他的男孩们开创性工作的致敬。
\end{quote}

\begin{quote}
\noindent\textbf{English.} The rediscovery of Mendel's work, and the breakthroughs that followed it, sparked a surge of interest in the social significance of genetics. While scientists had been grappling with the precise mechanisms of heredity through the eighteenth and nineteenth centuries, public concern had been mounting about the burden placed on society by what came to be called the "degenerate classes"---the inhabitants of poorhouses, workhouses, and insane asylums. What could be done with these people? It remained a matter of controversy whether they should be treated charitably---which, the less charitably inclined claimed, ensured such folk would never exert themselves and would therefore remain forever dependent on the largesse of the state or of private institutions---or whether they should be simply ignored, which, according to the charitably inclined, would result only in perpetuating the inability of the unfortunate to extricate themselves from their blighted circumstances.

\smallskip
\noindent\textbf{中文。} 孟德尔工作的重新发现以及随后的突破，引发了人们对遗传学社会意义的浓厚兴趣。十八、十九世纪期间，科学家一直在努力解决遗传的精确机制；与此同时，公众越来越担心所谓“退化阶层”给社会带来的负担，即济贫院、劳役所和精神病院中的居民。该如何对待这些人？一种主张是慈善救助；不那么有慈善心的人认为，这只会保证这些人永远不肯努力，从而永远依赖国家或私人机构的慷慨施舍。另一种主张是干脆不管；而有慈善心的人则认为，这只会让不幸者无法摆脱悲惨处境的状态永久化。争议一直存在。
\end{quote}

\begin{quote}
\noindent\textbf{English.} The publication of Darwin's \emph{Origin of Species} in 1859 brought these issues into sharp focus. Although Darwin carefully omitted to mention human evolution, fearing that to do so would only further inflame an already raging controversy, it required no great leap of imagination to apply his idea of natural selection to humans. Natural selection is the force that determines the fate of all genetic variations in nature---mutations like the one Morgan found in the fruit fly eye-color gene, but also perhaps differences in the abilities of human individuals to fend for themselves.

\smallskip
\noindent\textbf{中文。} 1859年达尔文《物种起源》的出版，使这些问题变得格外突出。达尔文担心提及人类进化只会进一步激化已经沸腾的争论，因此谨慎地避开了这一点；但把他的自然选择思想应用于人类，并不需要多大的想象力跳跃。自然选择是决定自然界一切遗传变异命运的力量，既包括摩尔根在果蝇眼色基因中发现的那类突变，也许也包括人类个体自谋生存能力上的差异。
\end{quote}

\begin{quote}
\noindent\textbf{English.} Natural populations have an enormous reproductive potential. Take fruit flies, with their generation time of just ten days, and females that produce some three hundred eggs apiece, half of which will be female: starting with a single fruit fly couple, after a month, that is, three generations later, you will have \(150 \times 150 \times 150\) fruit flies on your hands---that is more than 3 million flies, all of them derived from just one pair in just one month. Darwin made the point by choosing a species from the other end of the reproductive spectrum:

\smallskip
\noindent\textbf{中文。} 自然种群具有巨大的繁殖潜力。以果蝇为例，它们一代只需十天，每只雌蝇产卵约三百枚，其中一半会是雌性。从一对果蝇开始，一个月后，也就是三代以后，你手上会有 \(150 \times 150 \times 150\) 只果蝇，超过三百万只，而且全部在一个月内由最初一对繁衍而来。达尔文为了说明这一点，选择了繁殖谱另一端的物种：
\end{quote}

\begin{quote}
\noindent\textbf{English quotation.} The elephant is reckoned to be the slowest breeder of all known animals, and I have taken some pains to estimate its probable minimum rate of natural increase: it will be under the mark to assume that it breeds when thirty years old, and goes on breeding till ninety years old, bringing forth three pairs of young in this interval; if this be so, at the end of the fifth century there would be alive fifteen million elephants, descended from the first pair.

\smallskip
\noindent\textbf{中文引文。} 大象被认为是所有已知动物中繁殖最慢的；我曾费心估算它可能的最低自然增长率：假设它三十岁开始繁殖，并持续到九十岁，在此期间生下三对子代，这个估计还会偏低；若果真如此，到第五个世纪末，将会有一千五百万头大象存活，全都源自最初那一对。
\end{quote}

\begin{quote}
\noindent\textbf{English.} All these calculations assume that all the baby fruit flies and all the baby elephants make it successfully to adulthood. In theory, therefore, there must be an infinitely large supply of food and water to sustain this kind of reproductive overdrive. In reality, of course, those resources are limited, and not all baby fruit flies or baby elephants make it. There is competition among individuals within a species for those resources. What determines who wins the struggle for access to the resources? Darwin pointed out genetic variation means that some individuals have advantages in what he called "the struggle for existence." To take the famous example of Darwin's finches from the Galápagos Islands, those individuals with genetic advantages---like the right size of beak for eating the most abundant seeds---are more likely to survive and reproduce. So the advantageous genetic variant---having a bill the right size---tends to be passed on to the next generation. The result is that natural selection enriches the next generation with the beneficial mutation so that eventually, over enough generations, every member of the species ends up with that characteristic.

\smallskip
\noindent\textbf{中文。} 所有这些计算都假定，所有果蝇幼体和所有小象都能成功长到成年。因此在理论上，必须有无限多的食物和水来支撑这种繁殖狂飙。当然现实中这些资源是有限的，并不是所有果蝇幼体或小象都能活下来。同一物种的个体之间会为这些资源竞争。什么决定了谁在获取资源的斗争中获胜？达尔文指出，遗传变异意味着有些个体在他所谓的“生存斗争”中具有优势。以加拉帕戈斯群岛上著名的达尔文雀为例，具有遗传优势的个体，例如喙的大小正适合吃最丰富种子者，更有可能存活并繁殖。因此，有利的遗传变异，也就是拥有合适大小的喙，往往会传给下一代。结果是，自然选择让下一代中有益突变更加丰富；经过足够多代后，该物种每个成员最终都会拥有这一特征。
\end{quote}

\begin{quote}
\noindent\textbf{English.} The Victorians applied the same logic to humans. They looked around and were alarmed by what they saw. The decent, moral, hardworking middle classes were being massively out-reproduced by the dirty, immoral, lazy lower classes. The Victorians assumed that the virtues of decency, morality, and hard work ran in families just as the vices of filth, wantonness, and indolence did. Such characteristics must then be hereditary; thus, to the Victorians, morality and immorality were merely two of Darwin's genetic variants. And if the great unwashed were out-reproducing the respectable classes, then the "bad" genes would be increasing in the human population. The species was doomed! Humans would gradually become more and more depraved as the "immorality" gene became more and more common.

\smallskip
\noindent\textbf{中文。} 维多利亚时代的人把同样逻辑应用到人类身上。他们环顾四周，并对所见之事感到惊恐。体面、道德、勤劳的中产阶级，在繁殖数量上被肮脏、不道德、懒惰的下层阶级远远超过。维多利亚人假定，体面、道德和勤劳这些美德像污秽、放荡和懒惰这些恶习一样，会在家族中流传。那么这些特征必然是遗传的；于是对维多利亚人来说，道德和不道德不过是达尔文式遗传变异的两种。如果“不洗澡的大众”比体面阶级繁殖更多，那么“坏”基因就会在人类群体中增加。这个物种注定完了！随着“不道德”基因越来越常见，人类会逐渐变得越来越堕落。
\end{quote}

\begin{quote}
\noindent\textbf{English.} Francis Galton had good reason to pay special attention to Darwin's book, as the author was his cousin and friend. Darwin, some thirteen years older, had provided guidance during Galton's rather rocky college experience. But it was \emph{The Origin of Species} that would inspire Galton to start a social and genetic crusade that would ultimately have disastrous consequences. In 1883, a year after his cousin's death, Galton gave the movement a name: eugenics.

\smallskip
\noindent\textbf{中文。} 弗朗西斯·高尔顿有充分理由格外关注达尔文的书，因为作者既是他的表兄，也是他的朋友。年长约十三岁的达尔文曾在高尔顿颇为坎坷的大学经历中给予指导。但正是《物种起源》激励高尔顿发起了一场社会和遗传学十字军运动，其后果最终将十分灾难。1883年，也就是表兄去世一年后，高尔顿给这场运动取了一个名字：优生学。
\end{quote}

\begin{quote}
\noindent\textbf{English.} Eugenics was only one of Galton's many interests; Galton enthusiasts refer to him as a polymath, detractors as a dilettante. In fact, he made significant contributions to geography, anthropology, psychology, genetics, meteorology, statistics, and, by setting fingerprint analysis on a sound scientific footing, to criminology. Born in 1822 into a prosperous family, his education---partly in medicine and partly in mathematics---was mostly a chronicle of defeated expectations. The death of his father when he was twenty-one simultaneously freed him from paternal restraint and yielded a handsome inheritance; the young man duly took advantage of both. After a full six years of being what might be described today as a trust-fund dropout, however, Galton settled down to become a productive member of the Victorian establishment. He made his name leading an expedition to a then little-known region of southwest Africa in 1850--1852. In his account of his explorations, we encounter the first instance of the one strand that connects his many varied interests: he counted and measured everything. Galton was only happy when he could reduce a phenomenon to a set of numbers.

\smallskip
\noindent\textbf{中文。} 优生学只是高尔顿众多兴趣之一；崇拜他的人称他为博学多才者，批评者则称他为浅尝辄止的业余爱好者。事实上，他对地理学、人类学、心理学、遗传学、气象学、统计学都有重要贡献，并且通过把指纹分析建立在稳固的科学基础上，也贡献于犯罪学。高尔顿1822年出生于富裕家庭，他的教育部分是医学、部分是数学，却多半是一连串落空期望的记录。他二十一岁时父亲去世，这既使他摆脱了父亲约束，也给了他一笔可观遗产；年轻的高尔顿自然充分利用了二者。不过，在整整六年过着今天可称为“信托基金辍学者”的生活后，他安定下来，成为维多利亚体制中富有产出的一员。1850--1852年，他率领探险队前往当时鲜为人知的非洲西南部，并由此成名。在他对探险的记述中，我们首次看到贯穿他各种兴趣的一条线索：他对一切都计数和测量。只有当他能把一个现象化约为一组数字时，高尔顿才会高兴。
\end{quote}

\begin{quote}
\noindent\textbf{English.} At a missionary station he encountered a striking specimen of steatopygia---a condition of particularly protuberant buttocks, common among the indigenous Nama women of the region---and realized that this woman was naturally endowed with the figure that was then fashionable in Europe. The only difference was that it required enormous and costly ingenuity on the part of European dressmakers to create the desired "look" for their clients.

\smallskip
\noindent\textbf{中文。} 在一个传教站，他遇到了一位突出的脂臀体征者。脂臀是臀部特别隆起的状态，在当地土著纳马女性中很常见。高尔顿意识到，这位女性天生拥有当时欧洲流行的体形。唯一不同的是，欧洲裁缝若要为顾客制造这种理想“外观”，需要付出巨大而昂贵的巧思。
\end{quote}

\begin{quote}
\noindent\textbf{English quotation.} I profess to be a scientific man, and was exceedingly anxious to obtain accurate measurements of her shape; but there was a difficulty in doing this. I did not know a word of Hottentot [the Dutch name for the Nama], and could never therefore have explained to the lady what the object of my footrule could be; and I really dared not ask my worthy missionary host to interpret for me. I therefore felt in a dilemma as I gazed at her form, that gift of bounteous nature to this favoured race, which no mantua-maker, with all her crinoline and stuffing, can do otherwise than humbly imitate. The object of my admiration stood under a tree, and was turning herself about to all points of the compass, as ladies who wish to be admired usually do. Of a sudden my eye fell upon my sextant; the bright thought struck me, and I took a series of observations upon her figure in every direction, up and down, crossways, diagonally, and so forth, and I registered them carefully upon an outline drawing for fear of any mistake; this being done, I boldly pulled out my measuring tape, and measured the distance from where I was to the place she stood, and having thus obtained both base and angles, I worked out the results by trigonometry and logarithms.

\smallskip
\noindent\textbf{中文引文。} 我自称是科学人，极想取得她体形的准确测量；但这样做有困难。我一句霍屯督语（纳马人的荷兰称呼）也不会，因此无从向这位女士解释我的英尺尺究竟有什么目的；我也实在不敢请我可敬的传教士主人替我翻译。于是我望着她的形体陷入两难：那是慷慨自然赐予这个受惠种族的礼物，任何女装裁缝凭借裙撑和填充物，都只能谦卑模仿。我的赞赏对象站在一棵树下，像希望被欣赏的女士通常会做的那样，向罗盘的各个方向转身。突然我的目光落在六分仪上；一个绝妙念头击中了我，于是我从各个方向对她的身形做了一系列观测，上下、横向、斜向，等等，并小心地记录在轮廓图上，以免出错。做完这些，我大胆拿出卷尺，测量自己所在位置到她站立处的距离；这样取得底边和角度后，我便用三角学和对数算出了结果。
\end{quote}

\begin{figure}[htbp]
\centering
\ocrimage[width=0.60\linewidth]{images/ab6c7a6dc9c3960917faef512bf665aaa8b9de182031bc523ecdd38196dd229c.jpg}
\caption{A nineteenth-century exaggerated view of a Nama woman.\\十九世纪对一名纳马女性的夸张描绘。}
\end{figure}

\begin{quote}
\noindent\textbf{English.} Galton's passion for quantification resulted in his developing many of the fundamental principles of modern statistics. It also yielded some clever observations. For example, he tested the efficacy of prayer. He figured that if prayer worked, those most prayed for should be at an advantage; to test the hypothesis he studied the longevity of British monarchs. Every Sunday, congregations in the Church of England following the \emph{Book of Common Prayer} beseeched God to "Endue the king/queen plenteously with heavenly gifts; Grant him/her in health and wealth long to live." Surely, Galton reasoned, the cumulative effect of all those prayers should be beneficial. In fact, prayer seemed ineffectual: he found that on average the monarchs died somewhat younger than other members of the British aristocracy.

\smallskip
\noindent\textbf{中文。} 高尔顿对量化的热情，使他发展出现代统计学许多基本原则。它也产生了一些聪明观察。比如，他检验了祈祷的效力。他推想，如果祈祷有效，那么被祈祷最多的人应当有优势；为检验这一假说，他研究了英国君主的寿命。每个星期天，英国国教会众按照《公祷书》恳求上帝：“把丰盛的天赐礼物赋予国王/女王；赐予他/她健康与财富，使其长寿。”高尔顿推理道，所有这些祈祷的累积效果当然应该有益。事实上，祈祷似乎无效：他发现君主平均死亡年龄比英国贵族其他成员还略低。
\end{quote}

\begin{quote}
\noindent\textbf{English.} Because of the Darwin connection---their common grandfather, Erasmus Darwin, too, was one of the intellectual giants of his day---Galton was especially sensitive to the way in which certain lineages seemed to spawn disproportionately large numbers of prominent and successful people. In 1869 he published what would become the underpinning of all his ideas on eugenics, a treatise called \emph{Hereditary Genius: An Inquiry into Its Laws and Consequences}. In it he purported to show that talent, like simple genetic traits such as the Hapsburg Lip, does indeed run in families; he recounted, for example, how some families had produced generation after generation of judges. His analysis largely neglected to take into account the effect of the environment: the son of a prominent judge is, after all, rather more likely to become a judge---by virtue of his father's connections, if nothing else---than the son of a peasant farmer. Galton did not, however, completely overlook the effect of the environment, and it was he who first referred to the "nature/nurture" dichotomy, possibly in reference to Shakespeare's irredeemable villain, Caliban, "a devil, a born devil, on whose nature/Nurture can never stick."

\smallskip
\noindent\textbf{中文。} 由于与达尔文的关系，他们共同的祖父伊拉斯谟·达尔文也曾是当时代的思想巨人之一，高尔顿对某些家系似乎不成比例地产生大量杰出成功人物这一现象格外敏感。1869年，他发表了后来成为其全部优生学思想基础的著作《遗传的天才：其规律及后果研究》。在书中，他声称证明了才能确实像哈布斯堡唇这类简单遗传性状一样，在家族中流传；例如，他讲述某些家族如何一代又一代产生法官。他的分析在很大程度上忽略了环境影响：毕竟，一位显赫法官的儿子仅凭父亲的人脉，就比农民的儿子更可能成为法官。不过，高尔顿并没有完全忽略环境作用，也是他首先提出了“先天/后天”二分，也许是在回应莎士比亚笔下不可救赎的恶人凯列班：“一个魔鬼，天生的魔鬼，教养永远无法附着于其本性。”
\end{quote}

\begin{quote}
\noindent\textbf{English.} The results of his analysis, however, left no doubt in Galton's mind.

\smallskip
\noindent\textbf{中文。} 然而，他的分析结果在高尔顿心中没有留下任何疑问。
\end{quote}

\begin{quote}
\noindent\textbf{English quotation.} I have no patience with the hypothesis occasionally expressed, and often implied, especially in tales written to teach children to be good, that babies are born pretty much alike, and that the sole agencies in creating differences between boy and boy, and man and man, are steady application and moral effort. It is in the most unqualified manner that I object to pretensions of natural equality.

\smallskip
\noindent\textbf{中文引文。} 我不能容忍这样一种假说，它偶尔被明说，常常被暗示，尤其出现在教孩子做好人的故事中：婴儿出生时大体相同，而造成男孩与男孩、男人与男人之间差异的唯一因素，是持续努力和道德奋斗。我以最毫无保留的方式反对所谓自然平等的主张。
\end{quote}

\begin{quote}
\noindent\textbf{English.} A corollary of his conviction that these traits are genetically determined, he argued, was that it would be possible to "improve" the human stock by preferentially breeding gifted individuals, and preventing the less gifted from reproducing.

\smallskip
\noindent\textbf{中文。} 他认为，既然这些特征由基因决定，一个推论就是：通过优先让有天赋者繁殖，并阻止天赋较低者繁殖，就可能“改良”人类种群。
\end{quote}

\begin{quote}
\noindent\textbf{English quotation.} It is easy \(\ldots\) to obtain by careful selection a permanent breed of dogs or horses gifted with peculiar powers of running, or of doing anything else, so it would be quite practicable to produce a highly gifted race of men by judicious marriages during several consecutive generations.

\smallskip
\noindent\textbf{中文引文。} 通过仔细选择，很容易获得一个具有特殊奔跑能力或其他能力的稳定犬种或马种；因此，在连续若干代中通过明智婚配，培育出一个高度天赋的人类种族，也完全可行。
\end{quote}

\begin{quote}
\noindent\textbf{English.} Galton introduced the term eugenics, literally "good in birth," to describe this application of the basic principle of agricultural breeding to humans. In time, eugenics came to refer to "self-directed human evolution": by making conscious choices about who should have children, eugenicists believed that they could head off the "eugenic crisis" precipitated in the Victorian imagination by the high rates of reproduction of inferior stock coupled with the typically small families of the superior middle classes.

\smallskip
\noindent\textbf{中文。} 高尔顿引入“优生学”一词，字面意思是“出生良好”，用来描述把农业育种基本原则应用于人类。随着时间推移，优生学开始指“自我导向的人类进化”：通过有意识地选择谁应当生育，优生学者相信，他们能够阻止维多利亚想象中由劣等血统高繁殖率与优秀中产阶级家庭规模通常较小共同造成的“优生危机”。
\end{quote}

\begin{center}
\([\ldots]\)
\end{center}

\begin{figure}[htbp]
\centering
\ocrimage[width=0.72\linewidth]{images/5f5c2f9beaea0e59a190e44bd43fb9b305d148cb9108bddcf414f9951c876fb4.jpg}
\caption{Eugenics as it was perceived during the first part of the twentieth century: an opportunity for humans to control their own evolutionary destiny.\\二十世纪上半叶人们眼中的优生学：人类控制自身进化命运的机会。}
\end{figure}

\chapter{The Double Helix: This Is Life}
\begin{center}
\large 双螺旋：这就是生命
\end{center}

\begin{quote}
\noindent\textbf{English.} I got hooked on the gene during my third year at the University of Chicago. Until then, I had planned to be a naturalist and looked forward to a career far removed from the urban bustle of Chicago's South Side, where I grew up. My change of heart was inspired not by an unforgettable teacher but a little book that appeared in 1944, \emph{What Is Life?}, by the Austrian-born father of wave mechanics, Erwin Schrödinger. It grew out of several lectures he had given the year before at the Institute for Advanced Study in Dublin. That a great physicist had taken the time to write about biology caught my fancy. In those days, like most people, I considered chemistry and physics to be the "real" sciences, and theoretical physicists were science's top dogs.

\smallskip
\noindent\textbf{中文。} 我在芝加哥大学三年级时迷上了基因。在那之前，我原计划成为一名博物学家，期待一种远离芝加哥南区都市喧嚣的职业生涯；我正是在那里长大的。让我改变心意的不是一位令人难忘的老师，而是一本1944年出版的小书《生命是什么？》，作者是奥地利出生的波动力学之父埃尔温·薛定谔。此书源自他前一年在都柏林高等研究院所作的几场讲座。一位伟大的物理学家愿意花时间写生物学，这件事吸引了我。那时候，像大多数人一样，我认为化学和物理才是“真正的”科学，而理论物理学家则是科学界的顶尖人物。
\end{quote}

\begin{quote}
\noindent\textbf{English.} Schrödinger argued that life could be thought of in terms of storing and passing on biological information. Chromosomes were thus simply information bearers. Because so much information had to be packed into every cell, it must be compressed into what Schrödinger called a "hereditary code-script" embedded in the molecular fabric of chromosomes. To understand life, then, we would have to identify these molecules, and crack their code. He even speculated that understanding life---which would involve finding the gene---might take us beyond the laws of physics as we then understood them. Schrödinger's book was tremendously influential. Many of those who would become major players in Act 1 of molecular biology's great drama, including Francis Crick, a former physicist himself, had, like me, read \emph{What Is Life?} and been impressed.

\smallskip
\noindent\textbf{中文。} 薛定谔认为，可以从储存和传递生物信息的角度来理解生命。因此，染色体不过是信息载体。由于每个细胞必须容纳大量信息，这些信息必然被压缩进薛定谔所谓的“遗传密码脚本”中，并嵌入染色体的分子结构里。要理解生命，我们就必须识别这些分子，并破解它们的密码。他甚至推测，理解生命，也就是寻找基因，可能会把我们带出现有物理定律的范围。薛定谔的书影响极大。后来成为分子生物学伟大戏剧第一幕中主要人物的许多人，包括弗朗西斯·克里克这位前物理学家，都像我一样读过《生命是什么？》并深受触动。
\end{quote}

\begin{quote}
\noindent\textbf{English.} In my own case, Schrödinger struck a chord because I too was intrigued by the essence of life. A small minority of scientists still thought life depended upon a vital force emanating from an all-powerful god. But like most of my teachers, I disdained the very idea of vitalism. If such a "vital" force were calling the shots in nature's game, there was little hope life would ever be understood through the methods of science. On the other hand, the notion that life might be perpetuated by means of an instruction book inscribed in a secret code appealed to me. What sort of molecular code could be so elaborate as to convey all the multitudinous wonder of the living world? And what sort of molecular trick could ensure that the code is exactly copied every time a chromosome duplicates?

\smallskip
\noindent\textbf{中文。} 对我来说，薛定谔之所以拨动心弦，是因为我也被生命的本质所吸引。少数科学家仍认为，生命依赖一种来自全能上帝的生命力。但像我的多数老师一样，我轻视这种活力论观念。如果这样一种“生命”力量主宰自然游戏，那么通过科学方法理解生命就几乎没有希望。另一方面，生命也许依靠一本用秘密密码写成的说明书得以延续，这个想法吸引了我。什么样的分子密码能够复杂到传达生物世界如此繁多的奇妙？又是什么样的分子技巧，能够保证每次染色体复制时，密码都被精确复制？
\end{quote}

\begin{quote}
\noindent\textbf{English.} At the time of Schrödinger's Dublin lectures, most biologists supposed that proteins would eventually be identified as the primary bearers of genetic instruction. Proteins are molecular chains built up from twenty different building blocks, the amino acids. Because permutations in the order of amino acids along the chain are virtually infinite, proteins could, in principle, readily encode the information underpinning life's extraordinary diversity. DNA then was not considered a serious candidate for the bearer of code-scripts, even though it was exclusively located on chromosomes and had been known about for some seventy-five years. In 1869, Friedrich Miescher, a Swiss biochemist working in Germany, had isolated from pus-soaked bandages supplied by a local hospital a substance he called "nuclein." Because pus consists largely of white blood cells, which, unlike red blood cells, have nuclei and therefore DNA-containing chromosomes, Miescher had stumbled on a good source of DNA. When he later discovered that "nuclein" was to be found in chromosomes alone, Miescher understood that his discovery was indeed a big one. In 1893, he wrote: "Inheritance insures a continuity in form from generation to generation that lies even deeper than the chemical molecule. It lies in the structuring atomic groups. In this sense, I am a supporter of the chemical heredity theory."

\smallskip
\noindent\textbf{中文。} 在薛定谔都柏林讲座的时代，多数生物学家认为，蛋白质最终会被确认是遗传指令的主要承载者。蛋白质是由二十种不同构件，即氨基酸，组成的分子链。由于氨基酸在链上排列顺序的组合实际上无穷无尽，原则上蛋白质很容易编码支撑生命非凡多样性的信息。那时，DNA尚未被认真视为密码脚本的承载候选者，尽管它专门位于染色体上，并且已为人所知约七十五年。1869年，在德国工作的瑞士生物化学家弗里德里希·米歇尔，从当地医院提供的浸满脓液的绷带中分离出一种他称为“核素”的物质。脓液主要由白细胞构成；白细胞不同于红细胞，有细胞核，因此含有带DNA的染色体。米歇尔偶然找到了一个很好的DNA来源。后来他发现“核素”只存在于染色体中，便明白自己的发现确实意义重大。1893年，他写道：“遗传保证了形态从一代到下一代的连续性，这种连续性比化学分子更深；它存在于构成结构的原子团之中。从这个意义上说，我是化学遗传理论的支持者。”
\end{quote}

\begin{figure}[htbp]
\centering
\ocrimage[width=0.62\linewidth]{images/283af47520fa3108f7a0662a80b6c7e627dddf8a7e270d4eec97eebd5d0967f8.jpg}
\caption{The physicist Erwin Schrödinger, whose book \emph{What Is Life?} turned me on to the gene.\\物理学家埃尔温·薛定谔，他的《生命是什么？》让我迷上了基因。}
\end{figure}

\begin{quote}
\noindent\textbf{English.} Nevertheless, for decades afterward, chemistry would remain unequal to the task of analyzing the immense size and complexity of the DNA molecule. Only in the 1930s was DNA shown to be a long molecule containing four different chemical bases: adenine (A), guanine (G), thymine (T), and cytosine (C). But at the time of Schrödinger's lectures, it was still unclear just how the subunits, called deoxynucleotides, of the molecule were chemically linked. Nor was it known whether DNA molecules might vary in their sequences of the four different bases. If DNA were indeed Schrödinger's code-script, then the molecule would have to be capable of existing in an immense number of different forms. But back then it was still considered a possibility that one simple sequence like AGTC might be repeated over and over along the entire length of DNA chains.

\smallskip
\noindent\textbf{中文。} 尽管如此，此后数十年，化学仍无力分析DNA分子巨大的尺寸和复杂性。直到二十世纪三十年代，人们才证明DNA是一种长分子，含有四种不同化学碱基：腺嘌呤（A）、鸟嘌呤（G）、胸腺嘧啶（T）和胞嘧啶（C）。但在薛定谔讲座时，人们仍不清楚这种分子的亚单位，即脱氧核苷酸，究竟如何以化学方式连接。人们也不知道DNA分子在四种不同碱基的序列上是否可能变化。如果DNA确实是薛定谔所谓的密码脚本，那么这种分子就必须能够以数量巨大的不同形式存在。但当时仍有人认为，像AGTC这样简单的序列也许会沿DNA链全长一遍遍重复。
\end{quote}

\begin{quote}
\noindent\textbf{English.} DNA did not move into the genetic limelight until 1944, when Oswald Avery's lab at the Rockefeller Institute in New York City reported that the composition of the surface coats of pneumonia bacteria could be changed. This was not the result he and his junior colleagues, Colin MacLeod and Maclyn McCarty, expected.

\smallskip
\noindent\textbf{中文。} 直到1944年，DNA才进入遗传学聚光灯。当时纽约市洛克菲勒研究所奥斯瓦尔德·艾弗里的实验室报告说，肺炎细菌表面荚膜的组成可以被改变。这并不是他和年轻同事科林·麦克劳德、麦克林·麦卡蒂预期的结果。
\end{quote}

\begin{quote}
\noindent\textbf{English.} For more than a decade Avery's group had been following up on another most unexpected observation made in 1928 by Fred Griffith, a scientist in the British Ministry of Health. Griffith was interested in pneumonia and studied its bacterial agent, \emph{Pneumococcus}. It was known that there were two strains, designated "smooth" (S) and "rough" (R) according to their appearance under the microscope. These strains differed not only visually but also in their virulence. Inject S bacteria into a mouse, and within a few days the mouse dies; inject R bacteria and the mouse remains healthy. It turns out that S bacterial cells have a coating that prevents the mouse's immune system from recognizing the invader. The R cells have no such coating and are therefore readily attacked by the mouse's immune defenses.

\smallskip
\noindent\textbf{中文。} 十多年来，艾弗里团队一直在追踪英国卫生部科学家弗雷德·格里菲思1928年的另一个极其意外的观察。格里菲思关注肺炎，并研究其致病细菌肺炎球菌。人们已知这种细菌有两个菌株，根据显微镜下外观分别称为“光滑型”（S）和“粗糙型”（R）。这些菌株不仅外观不同，毒力也不同。把S型细菌注入小鼠，几天内小鼠死亡；注入R型细菌，小鼠仍保持健康。原来S型细胞有一层外衣，可以阻止小鼠免疫系统识别入侵者。R型细胞没有这层外衣，因此很容易被小鼠免疫防御攻击。
\end{quote}

\begin{figure}[htbp]
\centering
\ocrimage[width=0.70\linewidth]{images/0edbfbcce31b2e08b7233f24242c1aaad9c5441d31125ba58e18200cb469a047.jpg}
\caption{A view through the microscope of blood cells treated with a chemical that stains DNA. In order to maximize their oxygen-transporting capacity, red blood cells have no nucleus and therefore no DNA. But white blood cells, which patrol the bloodstream in search of intruders, have a nucleus containing chromosomes.\\显微镜下用DNA染色化学物质处理过的血细胞。为了最大限度提高输氧能力，红细胞没有细胞核，因此没有DNA。但在血流中巡逻、寻找入侵者的白细胞有细胞核，其中含有染色体。}
\end{figure}

\begin{quote}
\noindent\textbf{English.} Through his involvement with public health, Griffith knew that multiple strains had sometimes been isolated from a single patient, and so he was curious about how different strains might interact in his unfortunate mice. With one combination, he made a remarkable discovery: when he injected heat-killed S bacteria, harmless, and normal R bacteria, also harmless, the mouse died. How could two harmless forms of bacteria conspire to become lethal? The clue came when he isolated the \emph{Pneumococcus} bacteria retrieved from the dead mice and discovered living S bacteria. It appeared the living innocuous R bacteria had acquired something from the dead S variant; whatever it was, that something had allowed the R in the presence of the heat-killed S bacteria to transform itself into a living killer S strain. Griffith confirmed that this change was for real by culturing the S bacteria from the dead mouse over several generations: the bacteria bred true for the S type, just as any regular S strain would. A genetic change had indeed occurred to the R bacteria injected into the mouse.

\smallskip
\noindent\textbf{中文。} 由于参与公共卫生工作，格里菲思知道有时会从同一名患者体内分离出多个菌株，因此他很好奇不同菌株在他那些不幸的小鼠体内会如何相互作用。在一种组合中，他有了惊人发现：当他注入热杀死的S型细菌（无害）和正常R型细菌（也无害）时，小鼠却死了。两种无害细菌怎么会合谋致命？线索来自他从死亡小鼠体内分离肺炎球菌时发现了活的S型细菌。看来，活的无害R型细菌从死去的S型变体那里获得了某种东西；不管那是什么，它使R型在热杀死的S型细菌存在时，转化成了活的致命S型菌株。格里菲思通过把死鼠体内的S型细菌连续培养几代，确认这种改变是真实的：这些细菌像任何普通S型菌株一样稳定繁殖。注入小鼠的R型细菌确实发生了遗传改变。
\end{quote}

\begin{quote}
\noindent\textbf{English.} Though this transformation phenomenon seemed to defy all understanding, Griffith's observations at first created little stir in the scientific world. This was partly because Griffith was intensely private and so averse to large gatherings that he seldom attended scientific conferences. Once, he had to be virtually forced to give a lecture. Bundled into a taxi and escorted to the hall by colleagues, he discoursed in a mumbled monotone, emphasizing an obscure corner of his microbiological work but making no mention of bacterial transformation. Luckily, however, not everyone overlooked Griffith's breakthrough.

\smallskip
\noindent\textbf{中文。} 虽然这种转化现象看似完全不可理解，格里菲思的观察最初并未在科学界引起多大轰动。部分原因是格里菲思极其重视隐私，又非常厌恶大型聚会，几乎不参加科学会议。有一次，他几乎是被迫去做演讲的。同事把他塞进出租车并护送到会场后，他用含糊单调的声音讲述自己微生物学工作中一个冷僻角落，却丝毫没有提及细菌转化。不过幸运的是，并非所有人都忽视了格里菲思的突破。
\end{quote}

\begin{quote}
\noindent\textbf{English.} Oswald Avery was also interested in the sugarlike coats of the \emph{Pneumococcus}. He set out to duplicate Griffith's experiment in order to isolate and characterize whatever it was that had caused those R cells to change to the S type. In 1944 Avery, MacLeod, and McCarty published their results: an exquisite set of experiments showing unequivocally that DNA was the transforming principle. Culturing the bacteria in the test tube rather than in mice made it much easier to search for the chemical identity of the transforming factor in the heat-killed S cells. Methodically destroying one by one the biochemical components of the heat-treated S cells, Avery and his group looked to see whether transformation was prevented. First they degraded the sugarlike coat of the S bacteria. Transformation still occurred: the coat was not the transforming principle. Next they used a mixture of two protein-destroying enzymes, trypsin and chymotrypsin, to degrade virtually all the proteins in the S cells. To their surprise, transformation was again unaffected. Next they tried an enzyme, RNase, that breaks down RNA, ribonucleic acid, a second class of nucleic acids similar to DNA and possibly involved in protein synthesis. Again transformation occurred. Finally, they came to DNA, exposing the S bacterial extracts to the DNA-destroying enzyme, DNase. This time they hit a home run. All S-inducing activity ceased completely. The transforming factor was DNA.

\smallskip
\noindent\textbf{中文。} 奥斯瓦尔德·艾弗里也对肺炎球菌的糖样荚膜感兴趣。他开始重复格里菲思的实验，以分离并表征究竟是什么导致R型细胞变成S型。1944年，艾弗里、麦克劳德和麦卡蒂发表结果：一组精巧实验明确显示，DNA就是转化原则。与其在小鼠体内培养细菌，不如在试管中培养，这使他们更容易寻找热杀死S型细胞中转化因子的化学身份。艾弗里团队有条不紊地逐一破坏热处理S型细胞的生化成分，再观察转化是否被阻止。首先，他们降解S型细菌的糖样荚膜。转化仍然发生：荚膜不是转化原则。接着，他们使用两种破坏蛋白质的酶，胰蛋白酶和胰凝乳蛋白酶，几乎降解了S型细胞中所有蛋白质。令他们惊讶的是，转化仍不受影响。随后他们尝试RNase，这是一种分解RNA（核糖核酸）的酶；RNA是类似DNA的第二类核酸，可能参与蛋白质合成。转化又一次发生。最后，他们轮到DNA，把S型细菌提取物暴露给破坏DNA的酶DNase。这一次他们击中了靶心。所有诱导S型的活性完全停止。转化因子就是DNA。
\end{quote}

\begin{quote}
\noindent\textbf{English.} In part because of its bombshell implications, the resulting February 1944 paper by Avery, MacLeod, and McCarty met with a mixed response. Many geneticists accepted their conclusions. After all, DNA was found on every chromosome; why should it not be the genetic material? By contrast, however, most biochemists expressed doubt that DNA was a complex enough molecule to act as the repository of such a vast quantity of biological information. They continued to believe that proteins, the other component of chromosomes, would prove to be the hereditary substance. In principle, as the biochemists rightly noted, it would be much easier to encode a vast body of complex information using the twenty-letter amino-acid alphabet of proteins than the four-letter nucleotide alphabet of DNA. Particularly vitriolic in his rejection of DNA as the genetic substance was Avery's own colleague at the Rockefeller Institute, the protein chemist Alfred Mirsky. By then, however, Avery was no longer scientifically active. The Rockefeller Institute had mandatorily retired him at age sixty-five.

\smallskip
\noindent\textbf{中文。} 部分由于其含义像炸弹一样震撼，艾弗里、麦克劳德和麦卡蒂1944年2月发表的论文反响不一。许多遗传学家接受了他们的结论。毕竟，每条染色体上都有DNA；为什么它不能是遗传物质？相比之下，多数生物化学家怀疑DNA分子是否足够复杂，能作为如此巨大生物信息量的储存库。他们继续相信，染色体的另一种成分蛋白质才会被证明是遗传物质。原则上，正如生物化学家正确指出的，用蛋白质二十个字母的氨基酸字母表编码大量复杂信息，要比用DNA四个字母的核苷酸字母表容易得多。在否认DNA是遗传物质方面，尤其尖刻的是艾弗里在洛克菲勒研究所的同事、蛋白质化学家阿尔弗雷德·米尔斯基。然而，到那时艾弗里已经不再活跃于科学研究。洛克菲勒研究所强制他在六十五岁退休。
\end{quote}

\begin{quote}
\noindent\textbf{English.} Avery missed out on more than the opportunity to defend his work against the attacks of his colleagues: he was never awarded the Nobel Prize, which was certainly his due, for identifying DNA as the transforming principle. Because the Nobel committee makes its records public fifty years following each award, we now know that Avery's candidacy was blocked by the Swedish physical chemist Einar Hammarsten. Though Hammarsten's reputation was based largely on his having produced DNA samples of unprecedented high quality, he still believed genes to be an undiscovered class of proteins. In fact, even after the double helix was found, Hammarsten continued to insist that Avery should not receive the prize until after the mechanism of DNA transformation had been completely worked out. Avery died in 1955; had he lived only a few more years, he would almost certainly have gotten the prize.

\smallskip
\noindent\textbf{中文。} 艾弗里错过的不只是为自己的工作抵御同事攻击的机会。由于识别出DNA是转化原则，他理应获得诺贝尔奖，却从未获奖。诺贝尔委员会会在每次授奖五十年后公开档案，因此我们现在知道，艾弗里的候选资格被瑞典物理化学家埃纳尔·哈马斯滕阻止了。哈马斯滕的声誉主要建立在他制备出前所未有高质量DNA样品之上，但他仍相信基因是一类尚未发现的蛋白质。事实上，即便双螺旋被发现以后，哈马斯滕仍坚持认为，在DNA转化机制被完全弄清楚之前，不应把奖授予艾弗里。艾弗里1955年去世；如果他再多活几年，几乎肯定会获奖。
\end{quote}

\begin{quote}
\noindent\textbf{English.} When I arrived at Indiana University in the fall of 1947 with plans to pursue the gene for my Ph.D. thesis, Avery's paper came up over and over in conversations. By then, no one doubted the reproducibility of his results, and more recent work coming out of the Rockefeller Institute made it all the less likely that proteins would prove to be the genetic actors in bacterial transformation. DNA had at last become an important objective for chemists setting their sights on the next breakthrough. In Cambridge, England, the canny Scottish chemist Alexander Todd rose to the challenge of identifying the chemical bonds that linked together nucleotides in DNA. By early 1951, his lab had proved that these links were always the same, such that the backbone of the DNA molecule was very regular. During the same period, the Austrian-born refugee Erwin Chargaff, at the College of Physicians and Surgeons of Columbia University, used the new technique of paper chromatography to measure the relative amounts of the four DNA bases in DNA samples extracted from a variety of vertebrates and bacteria. While some species had DNA in which adenine and thymine predominated, others had DNA with more guanine and cytosine. The possibility thus presented itself that no two DNA molecules had the same composition.

\smallskip
\noindent\textbf{中文。} 1947年秋天，我来到印第安纳大学，计划以基因为博士论文主题，艾弗里的论文在谈话中一次又一次被提起。到那时，已经没有人怀疑他的结果可以重复；洛克菲勒研究所新近的工作也使蛋白质会被证明是细菌转化中遗传行动者的可能性越来越小。DNA终于成为化学家瞄准下一次突破时的重要目标。在英国剑桥，精明的苏格兰化学家亚历山大·托德迎接了挑战：识别把DNA中核苷酸连接起来的化学键。到1951年初，他的实验室已经证明这些连接始终相同，因此DNA分子的骨架非常规则。同一时期，出生于奥地利的难民埃尔温·查伽夫在哥伦比亚大学内外科医学院，利用纸层析新技术，测量从多种脊椎动物和细菌中提取的DNA样品里四种DNA碱基的相对含量。有些物种的DNA中腺嘌呤和胸腺嘧啶占优势，另一些物种的DNA则含有更多鸟嘌呤和胞嘧啶。由此出现了一种可能：没有两个DNA分子的组成完全相同。
\end{quote}

\begin{quote}
\noindent\textbf{English.} At Indiana I joined a small group of visionary scientists, mostly physicists and chemists, studying the reproductive process of the viruses that attack bacteria, bacteriophages---"phages" for short. The Phage Group was born when my Ph.D. supervisor, the Italian-trained medic Salvador Luria, and his close friend, the German-born theoretical physicist Max Delbrück, teamed up with the American physical chemist Alfred Hershey. During World War II both Luria and Delbrück were considered enemy aliens, and thus ineligible to serve in the war effort of American science, even though Luria, a Jew, had been forced to leave France for New York City and Delbrück had fled Germany as an objector to Nazism. Thus excluded, they continued to work in their respective university labs---Luria at Indiana and Delbrück at Vanderbilt---and collaborated on phage experiments during successive summers at Cold Spring Harbor. In 1943, they joined forces with the brilliant but taciturn Hershey, then doing phage research of his own at Washington University in St. Louis.

\smallskip
\noindent\textbf{中文。} 在印第安纳，我加入了一小群有远见的科学家，他们多半是物理学家和化学家，研究攻击细菌的病毒，即噬菌体，简称“噬菌体”的繁殖过程。噬菌体小组诞生于我的博士导师、受意大利训练的医生萨尔瓦多·卢里亚，与他的密友、德国出生的理论物理学家马克斯·德尔布吕克，以及美国物理化学家阿尔弗雷德·赫尔希联合之时。第二次世界大战期间，卢里亚和德尔布吕克都被视为敌国侨民，因此没有资格服务于美国科学的战时努力，尽管身为犹太人的卢里亚被迫从法国逃往纽约，德尔布吕克则因反对纳粹主义而逃离德国。于是，被排除在外的他们继续在各自大学实验室工作，卢里亚在印第安纳，德尔布吕克在范德比尔特，并在连续几个夏天于冷泉港合作开展噬菌体实验。1943年，他们与才华横溢但沉默寡言的赫尔希联手；赫尔希当时正在圣路易斯华盛顿大学从事自己的噬菌体研究。
\end{quote}

\begin{quote}
\noindent\textbf{English.} The Phage Group's program was based on its belief that phages, like all viruses, were in effect naked genes. This concept had first been proposed in 1922 by the imaginative American geneticist Herman J. Muller, who three years later demonstrated that X rays cause mutations. His belated Nobel Prize came in 1946, just after he joined the faculty of Indiana University. It was his presence, in fact, that led me to Indiana. Having started his career under T. H. Morgan, Muller knew better than anyone else how genetics had evolved during the first half of the twentieth century, and I was enthralled by his lectures during my first term. His work on fruit flies, \emph{Drosophila}, however, seemed to me to belong more to the past than to the future, and I only briefly considered doing thesis research under his supervision. I opted instead for Luria's phages, an even speedier experimental subject than \emph{Drosophila}: genetic crosses of phages done one day could be analyzed the next.

\smallskip
\noindent\textbf{中文。} 噬菌体小组的计划建立在这样一种信念上：噬菌体像所有病毒一样，实际上是裸露的基因。这个概念最早由富有想象力的美国遗传学家赫尔曼·J. 穆勒在1922年提出；三年后，他证明X射线会导致突变。他迟来的诺贝尔奖在1946年到来，就在他加入印第安纳大学教职之后。事实上，正是他的存在把我带到了印第安纳。穆勒在T. H. 摩尔根门下开始职业生涯，他比任何人都更清楚遗传学在二十世纪上半叶如何发展；我第一学期听他的课时深深着迷。不过，他关于果蝇的工作在我看来更属于过去而不是未来，我只是短暂考虑过在他指导下做论文研究。我最终选择了卢里亚的噬菌体，它比果蝇还是更快的实验对象：一天完成的噬菌体遗传杂交，第二天就能分析。
\end{quote}

\begin{quote}
\noindent\textbf{English.} For my Ph.D. thesis research, Luria had me follow in his footsteps by studying how X rays killed phage particles. Initially I had hoped to show that viral death was caused by damage to phage DNA. Reluctantly, however, I eventually had to concede that my experimental approach could never give unambiguous answers at the chemical level. I could draw only biological conclusions. Even though phages were indeed effectively naked genes, I realized that the deep answers the Phage Group was seeking could be arrived at only through advanced chemistry. DNA somehow had to transcend its status as an acronym; it had to be understood as a molecular structure in all its chemical detail.

\smallskip
\noindent\textbf{中文。} 我的博士论文研究中，卢里亚让我沿着他的道路，研究X射线如何杀死噬菌体颗粒。起初我希望证明，病毒死亡是由噬菌体DNA受损造成的。然而，我最终不得不勉强承认，我的实验方法永远无法在化学层面给出明确答案。我只能得出生物学结论。尽管噬菌体确实等同于裸露基因，我意识到，噬菌体小组寻求的深层答案只能通过高级化学获得。DNA必须以某种方式超越其缩写词地位；它必须被理解为一个具有全部化学细节的分子结构。
\end{quote}

\begin{quote}
\noindent\textbf{English.} Upon finishing my thesis, I saw no alternative but to move to a lab where I could study DNA chemistry. Unfortunately, however, knowing almost no pure chemistry, I would have been out of my depth in any lab attempting difficult experiments in organic or physical chemistry. I therefore took a postdoctoral fellowship in the Copenhagen lab of the biochemist Herman Kalckar in the fall of 1950. He was studying the synthesis of the small molecules that make up DNA, but I figured out quickly that his biochemical approach would never lead to an understanding of the essence of the gene. Every day spent in his lab would be one more day's delay in learning how DNA carried genetic information.

\smallskip
\noindent\textbf{中文。} 完成论文后，我看不到别的选择，只能转到一个可以研究DNA化学的实验室。不幸的是，我几乎不懂纯化学，因此在任何尝试进行艰难有机化学或物理化学实验的实验室里都会力不从心。于是1950年秋，我接受了生物化学家赫尔曼·卡尔卡在哥本哈根实验室的博士后职位。他研究构成DNA的小分子的合成，但我很快看出，他的生化方法永远无法引向对基因本质的理解。在他的实验室里每多待一天，就会让学习DNA如何携带遗传信息的时间再推迟一天。
\end{quote}

\begin{quote}
\noindent\textbf{English.} My Copenhagen year nonetheless ended productively. To escape the cold Danish spring, I went to the Zoological Station at Naples during April and May. During my last week there, I attended a small conference on X-ray diffraction methods for determining the 3-D structure of molecules. X-ray diffraction is a way of studying the atomic structure of any molecule that can be crystallized. The crystal is bombarded with X rays, which bounce off its atoms and are scattered. The scatter pattern gives information about the structure of the molecule but, taken alone, is not enough to solve the structure. The additional information needed is the "phase assignment," which deals with the wave properties of the molecule. Solving the phase problem was not easy, and at that time only the most audacious scientists were willing to take it on. Most of the successes of the diffraction method had been achieved with relatively simple molecules.

\smallskip
\noindent\textbf{中文。} 尽管如此，我在哥本哈根的那一年还是以富有成效的方式结束了。为了逃离丹麦寒冷的春天，我在四月和五月去了那不勒斯动物学站。在那里的最后一周，我参加了一个关于用X射线衍射方法测定分子三维结构的小型会议。X射线衍射是一种研究任何可结晶分子原子结构的方法。晶体受到X射线轰击，X射线从其原子上反弹并散射。散射图案提供分子结构信息，但单凭它还不足以解出结构。还需要额外信息，即“相位指定”，它涉及分子的波动性质。解决相位问题并不容易，在当时只有最大胆的科学家愿意承担。衍射法的大多数成功，都发生在相对简单的分子上。
\end{quote}

\begin{quote}
\noindent\textbf{English.} My expectations for the conference were low. I believed that a three-dimensional understanding of protein structure, or for that matter of DNA, was more than a decade away. Disappointing earlier X-ray photos suggested that DNA was particularly unlikely to yield up its secrets via the X-ray approach. These results were not surprising since the exact sequences of DNA were expected to differ from one individual molecule to another. The resulting irregularity of surface configurations would understandably prevent the long thin DNA chains from lying neatly side by side in the regular repeating patterns required for X-ray analysis to be successful.

\smallskip
\noindent\textbf{中文。} 我对这次会议的期待很低。我认为，对蛋白质结构的三维理解，甚至对DNA的三维理解，至少还要十多年。早先令人失望的X射线照片提示，DNA尤其不可能通过X射线方法交出秘密。这些结果并不意外，因为人们预期不同DNA分子的精确序列会彼此不同。由此造成的表面构型不规则，可以理解地会阻止细长DNA链整齐并排排列成X射线分析成功所需的规则重复图案。
\end{quote}

\begin{quote}
\noindent\textbf{English.} It was therefore a surprise and a delight to hear the last-minute talk on DNA by a thirty-four-year-old Englishman named Maurice Wilkins from the Biophysics Lab of King's College, London. Wilkins was a physicist who during the war had worked on the Manhattan Project. For him, as for many of the other scientists involved, the actual deployment of the bomb on Hiroshima and Nagasaki, supposedly the culmination of all their work, was profoundly disillusioning.

\smallskip
\noindent\textbf{中文。} 因此，听到伦敦国王学院生物物理实验室三十四岁的英国人莫里斯·威尔金斯临时作关于DNA的报告，既出乎意料又令人高兴。威尔金斯是物理学家，战争期间曾参与曼哈顿计划。对他来说，如同对许多参与其中的科学家一样，原本被视为全部工作巅峰的炸弹实际投向广岛和长崎，令人深感幻灭。
\end{quote}

\begin{figure}[htbp]
\centering
\ocrimage[width=0.66\linewidth]{images/a27a6630519567029c486e0a8709a2111b58555b107bf1f6aa5c6e55aa3f1623.jpg}
\caption{Maurice Wilkins in his lab at King's College, London.\\莫里斯·威尔金斯在伦敦国王学院的实验室。}
\end{figure}

\begin{quote}
\noindent\textbf{English.} He considered forsaking science altogether to become a painter in Paris, but biology intervened. He too had read Schrödinger's book, and was now tackling DNA with X-ray diffraction. He displayed a photograph of an X-ray diffraction pattern he had recently obtained, and its many precise reflections indicated a highly regular crystalline packing. DNA, one had to conclude, must have a regular structure, the elucidation of which might well reveal the nature of the gene. Instantly I saw myself moving to London to help Wilkins find the structure. My attempts to converse with him after his talk, however, went nowhere. All I got for my efforts was a declaration of his conviction that much hard work lay ahead.

\smallskip
\noindent\textbf{中文。} 他曾考虑彻底放弃科学，去巴黎当画家，但生物学介入了。他也读过薛定谔的书，如今正用X射线衍射研究DNA。他展示了一张自己最近获得的X射线衍射图，许多精确反射表明其中存在高度规则的晶体堆积。人们不得不下结论：DNA一定有规则结构，而阐明这种结构很可能揭示基因的本质。我立刻想象自己搬到伦敦，帮助威尔金斯寻找结构。然而，他报告后我试图与他交谈，却毫无进展。我努力所得，只是他坚称前方还有大量艰苦工作。
\end{quote}

\begin{quote}
\noindent\textbf{English.} While I was hitting consecutive dead ends, back in America the world's preeminent chemist, Caltech's Linus Pauling, announced a major triumph: he had found the exact arrangement in which chains of amino acids, called polypeptides, fold up in proteins, and called his structure the \(\alpha\)-helix. That it was Pauling who made this breakthrough was no surprise: he was a scientific superstar. His book \emph{The Nature of the Chemical Bond} essentially laid the foundation of modern chemistry, and, for chemists of the day, it was the Bible. Pauling had been a precocious child. When he was nine, his father, a druggist in Oregon, wrote to the \emph{Oregonian} newspaper requesting suggestions of reading matter for his bookish son, adding that he had already read the Bible and Darwin's \emph{Origin of Species}. But the early death of Pauling's father, which brought the family to financial ruin, makes it remarkable that the promising young man managed to get an education at all.

\smallskip
\noindent\textbf{中文。} 当我接连碰壁时，在美国，世界最杰出的化学家、加州理工学院的莱纳斯·鲍林宣布了一项重大胜利：他找到了氨基酸链，也称多肽链，在蛋白质中折叠的确切排列方式，并把这一结构称为 \(\alpha\)-螺旋。取得这一突破的人是鲍林并不令人意外：他是科学巨星。他的《化学键的本质》基本奠定了现代化学基础，对当时的化学家来说就是圣经。鲍林童年早慧。九岁时，他在俄勒冈当药剂师的父亲写信给《俄勒冈人报》，请他们为爱读书的儿子推荐读物，并补充说孩子已经读过《圣经》和达尔文的《物种起源》。但鲍林父亲早逝，使家庭陷入经济破产，这也让这个前途光明的年轻人竟能接受教育一事显得格外不易。
\end{quote}

\begin{quote}
\noindent\textbf{English.} As soon as I returned to Copenhagen I read about Pauling's \(\alpha\)-helix. To my surprise, his model was not based on a deductive leap from experimental X-ray diffraction data. Instead, it was Pauling's long experience as a structural chemist that had emboldened him to infer which type of helical fold would be most compatible with the underlying chemical features of the polypeptide chain. Pauling made scale models of the different parts of the protein molecule, working out plausible schemes in three dimensions. He had reduced the problem to a kind of three-dimensional jigsaw puzzle in a way that was simple yet brilliant.

\smallskip
\noindent\textbf{中文。} 我一回到哥本哈根，就读到了鲍林的 \(\alpha\)-螺旋。令我惊讶的是，他的模型并不是基于从X射线衍射实验数据做出的演绎飞跃。相反，是鲍林作为结构化学家的长期经验，使他敢于推断哪一种螺旋折叠最能与多肽链底层化学特征相容。鲍林制作了蛋白质分子不同部分的比例模型，在三维空间中推敲可行方案。他把问题化约成一种三维拼图，方式简单却高明。
\end{quote}

\begin{quote}
\noindent\textbf{English.} Whether the \(\alpha\)-helix was correct---in addition to being pretty---was now the question. Only a week later, I got the answer. Sir Lawrence Bragg, the English inventor of X-ray crystallography and 1915 Nobel laureate in Physics, came to Copenhagen and excitedly reported that his junior colleague, the Austrian-born chemist Max Perutz, had ingeniously used synthetic polypeptides to confirm the correctness of Pauling's \(\alpha\)-helix. It was a bittersweet triumph for Bragg's Cavendish Laboratory. The year before, they had completely missed the boat in their paper outlining possible helical folds for polypeptide chains.

\smallskip
\noindent\textbf{中文。} \(\alpha\)-螺旋除了漂亮之外是否正确，成了问题。仅一周后，我得到了答案。X射线晶体学的英国发明者、1915年诺贝尔物理学奖得主劳伦斯·布拉格爵士来到哥本哈根，兴奋地报告说，他的年轻同事、奥地利出生的化学家马克斯·佩鲁茨巧妙地利用合成多肽，证实了鲍林 \(\alpha\)-螺旋的正确性。对布拉格的卡文迪许实验室来说，这是一次苦乐参半的胜利。就在前一年，他们在一篇概述多肽链可能螺旋折叠的论文中完全错失了良机。
\end{quote}

\begin{figure}[htbp]
\centering
\ocrimage[width=0.70\linewidth]{images/214fa4b29d64e67c2fdb9a065099f912ba77e0f50d9c2c3553ba6b0c32ff357d.jpg}
\caption{Lawrence Bragg (left) with Linus Pauling, who is carrying a model of the \(\alpha\)-helix.\\劳伦斯·布拉格（左）与莱纳斯·鲍林；鲍林手中拿着 \(\alpha\)-螺旋模型。}
\end{figure}

\begin{quote}
\noindent\textbf{English.} By then Salvador Luria had tentatively arranged for me to take up a research position at the Cavendish. Located at Cambridge University, this was the most famous laboratory in all of science. Here Ernest Rutherford first described the structure of the atom. Now it was Bragg's own domain, and I was to work as apprentice to the English chemist John Kendrew, who was interested in determining the 3-D structure of the protein myoglobin. Luria advised me to visit the Cavendish as soon as possible. With Kendrew in the States, Max Perutz would check me out. Together, Kendrew and Perutz had earlier established the Medical Research Council (MRC) Unit for the Study of the Structure of Biological Systems.

\smallskip
\noindent\textbf{中文。} 到那时，萨尔瓦多·卢里亚已初步安排我去卡文迪许实验室任研究职位。卡文迪许位于剑桥大学，是全科学界最著名的实验室。欧内斯特·卢瑟福最早就是在这里描述原子结构。如今这里是布拉格的地盘，而我将作为英国化学家约翰·肯德鲁的学徒工作；肯德鲁关心的是确定肌红蛋白的三维结构。卢里亚建议我尽快访问卡文迪许。肯德鲁在美国，所以马克斯·佩鲁茨会来考察我。肯德鲁和佩鲁茨早些时候共同建立了医学研究委员会（MRC）生物系统结构研究单元。
\end{quote}

\begin{quote}
\noindent\textbf{English.} A month later in Cambridge, Perutz assured me that I could quickly master the necessary X-ray diffraction theory and should have no difficulty fitting in with the others in their tiny MRC Unit. To my relief, he was not put off by my biology background. Nor was Lawrence Bragg, who briefly came down from his office to look me over.

\smallskip
\noindent\textbf{中文。} 一个月后在剑桥，佩鲁茨向我保证，我可以很快掌握必要的X射线衍射理论，并且在他们那个小小的MRC单元里融入其他人不会有困难。令我松了一口气的是，他并没有被我的生物学背景吓退。劳伦斯·布拉格也没有；他短暂地从办公室下来打量了我一番。
\end{quote}

\begin{quote}
\noindent\textbf{English.} I was twenty-three when I arrived back at the MRC Unit in Cambridge in early October. I found myself sharing space in the biochemistry room with a thirty-five-year-old ex-physicist, Francis Crick, who had spent the war working on magnetic mines for the Admiralty. When the war ended, Crick had planned to stay on in military research, but, on reading Schrödinger's \emph{What Is Life?}, he had moved toward biology. Now he was at the Cavendish to pursue the 3-D structure of proteins for his Ph.D.

\smallskip
\noindent\textbf{中文。} 十月初我回到剑桥MRC单元时二十三岁。我发现自己在生化室里与一位三十五岁的前物理学家弗朗西斯·克里克共用空间；战争期间，他为海军部研究磁性水雷。战争结束时，克里克本打算继续从事军事研究，但读了薛定谔的《生命是什么？》后，他转向了生物学。如今他在卡文迪许攻读博士，研究蛋白质三维结构。
\end{quote}

\begin{quote}
\noindent\textbf{English.} Crick was always fascinated by the intricacies of important problems. His endless questions as a child compelled his weary parents to buy him a children's encyclopedia, hoping that it would satisfy his curiosity. But it only made him insecure: he confided to his mother his fear that everything would have been discovered by the time he grew up, leaving him nothing to do. His mother reassured him, correctly, as it happened, that there would still be a thing or two for him to figure out.

\smallskip
\noindent\textbf{中文。} 克里克总是被重要问题的复杂细节所吸引。小时候，他无穷无尽的问题迫使疲惫的父母给他买了一套儿童百科全书，希望满足他的好奇心。但这只让他更不安：他向母亲吐露担忧，害怕自己长大时一切都已被发现，留给他无事可做。他母亲安慰他说，仍会有一两件事等着他弄明白；事实证明她是对的。
\end{quote}

\begin{quote}
\noindent\textbf{English.} A great talker, Crick was invariably the center of attention in any gathering. His booming laugh was forever echoing down the hallways of the Cavendish. As the MRC Unit's resident theoretician, he used to come up with a novel insight at least once a month, and he would explain his latest idea at great length to anyone willing to listen. The morning we met he lit up when he learned that my objective in coming to Cambridge was to learn enough crystallography to have a go at the DNA structure. Soon I was asking Crick's opinion about using Pauling's model-building approach to go directly for the structure. Would we need many more years of diffraction experimentation before modeling would be practicable? To bring us up to speed on the status of DNA structural studies, Crick invited Maurice Wilkins, a friend since the end of the war, up from London for Sunday lunch. Then we could learn what progress Wilkins had made since his talk in Naples.

\smallskip
\noindent\textbf{中文。} 克里克极其健谈，在任何聚会中总是注意力中心。他洪亮的笑声永远回荡在卡文迪许走廊中。作为MRC单元的常驻理论家，他过去每月至少会想出一个新见解，并向任何愿意听的人详尽解释自己的最新想法。我们相遇的那个早晨，当他得知我来剑桥的目标是学习足够的晶体学，以便尝试DNA结构时，他立刻来了精神。很快，我就询问克里克是否可以采用鲍林的模型构建方法，直接奔向结构。在建模可行之前，我们还需要许多年衍射实验吗？为了让我们了解DNA结构研究的现状，克里克邀请战后以来的朋友莫里斯·威尔金斯从伦敦来吃星期日午餐。这样我们就可以知道威尔金斯自那不勒斯报告以来取得了什么进展。
\end{quote}

\begin{figure}[htbp]
\centering
\ocrimage[width=0.64\linewidth]{images/3544e5837f922455f23bb7338a782d2a404c84d9e43c0ad947e0ab5094d455f8.jpg}
\caption{Francis Crick with the Cavendish X-ray tube.\\弗朗西斯·克里克与卡文迪许的X射线管。}
\end{figure}

\begin{quote}
\noindent\textbf{English.} Wilkins expressed his belief that DNA's structure was a helix, formed by several chains of linked nucleotides twisted around each other. All that remained to be settled was the number of chains. At the time, Wilkins favored three on the basis of his density measurements of DNA fibers. He was keen to start model-building, but he had run into a roadblock in the form of a new addition to the King's College Biophysics Unit, Rosalind Franklin.

\smallskip
\noindent\textbf{中文。} 威尔金斯表示，他相信DNA结构是螺旋，由几条相连核苷酸链相互缠绕形成。剩下需要确定的只是链的数量。当时，基于对DNA纤维的密度测量，威尔金斯倾向于三链。他急于开始建模，但遇到了一道障碍：国王学院生物物理单元新来的人，罗莎琳德·富兰克林。
\end{quote}

\begin{quote}
\noindent\textbf{English.} A thirty-one-year-old Cambridge-trained physical chemist, Franklin was an obsessively professional scientist; for her twenty-ninth birthday all she requested was her own subscription to her field's technical journal, \emph{Acta Crystallographica}. Logical and precise, she was impatient with those who acted otherwise. And she was given to strong opinions, once describing her Ph.D. thesis adviser, Ronald Norrish, a future Nobel laureate, as "stupid, bigoted, deceitful, ill-mannered and tyrannical." Outside the laboratory, she was a determined and gutsy mountaineer, and, coming from the upper echelons of London society, she belonged to a more rarefied social world than most scientists. At the end of a hard day at the bench, she would occasionally change out of her lab coat into an elegant evening gown and disappear into the night.

\smallskip
\noindent\textbf{中文。} 富兰克林三十一岁，是受剑桥训练的物理化学家，也是一位近乎执着的职业科学家；她二十九岁生日时唯一要求的礼物，是订阅自己领域的技术期刊《晶体学报》。她逻辑严密、精确，对行事不如此的人缺乏耐心。她观点强烈，曾把自己的博士论文导师、后来的诺贝尔奖得主罗纳德·诺里什形容为“愚蠢、偏执、虚伪、粗鲁而专横”。在实验室之外，她是坚定而勇敢的登山者；出身伦敦上流社会的她，属于比多数科学家更精致的社交世界。辛苦工作一天后，她偶尔会脱下实验服，换上优雅晚礼服，消失在夜色中。
\end{quote}

\begin{figure}[htbp]
\centering
\ocrimage[width=0.66\linewidth]{images/10e89428031ee62695360770c2642e8814912404728b22a692edeafa9def4bed.jpg}
\caption{Rosalind Franklin on one of the mountain hiking vacations she loved.\\罗莎琳德·富兰克林在她热爱的山地徒步假期中。}
\end{figure}

\begin{quote}
\noindent\textbf{English.} Just back from a four-year X-ray crystallographic investigation of graphite in Paris, Franklin had been assigned to the DNA project while Wilkins was away from King's. Unfortunately, the pair soon proved incompatible. Franklin, direct and data-focused, and Wilkins, retiring and speculative, were destined never to collaborate. Shortly before Wilkins accepted our lunch invitation, the two had had a big blowup in which Franklin had insisted that no model-building could commence before she collected much more extensive diffraction data. Now they effectively did not communicate, and Wilkins would have no chance to learn of her progress until Franklin presented her lab seminar scheduled for the beginning of November. If we wanted to listen, Crick and I were welcome to go as Wilkins's guests.

\smallskip
\noindent\textbf{中文。} 富兰克林刚从巴黎回来，完成了为期四年的石墨X射线晶体学研究；在威尔金斯离开国王学院期间，她被分配到DNA项目。不幸的是，两人很快证明彼此不合。富兰克林直截了当、专注数据；威尔金斯内向、倾向推测，他们注定无法合作。就在威尔金斯接受我们的午餐邀请前不久，两人大吵了一架；富兰克林坚持认为，在她收集到更多广泛的衍射数据之前，不能开始任何模型构建。现在他们实际上不再交流，而威尔金斯也无从了解她的进展，直到富兰克林按计划在十一月初作实验室报告。如果我们想听，克里克和我可以作为威尔金斯的客人去。
\end{quote}

\begin{quote}
\noindent\textbf{English.} Crick was unable to make the seminar, so I attended alone and briefed him later on what I believed to be its key take-home messages on crystalline DNA. In particular, I described from memory Franklin's measurements of the crystallographic repeats and the water content. This prompted Crick to begin sketching helical grids on a sheet of paper, explaining that the new helical X-ray theory he had devised with Bill Cochran and Vladimir Vand would permit even me, a former bird-watcher, to predict correctly the diffraction patterns expected from the molecular models we would soon be building at the Cavendish.

\smallskip
\noindent\textbf{中文。} 克里克无法参加报告，所以我独自前往，后来向他转述我认为其中关于晶态DNA的关键要点。尤其是，我凭记忆描述了富兰克林对晶体学重复周期和含水量的测量。这促使克里克在纸上开始画螺旋网格，并解释说，他与比尔·科克伦和弗拉基米尔·范德共同提出的新螺旋X射线理论，甚至能让我这个前观鸟者，正确预测我们很快将在卡文迪许搭建的分子模型应产生的衍射图案。
\end{quote}

\begin{quote}
\noindent\textbf{English.} As soon as we got back to Cambridge, I arranged for the Cavendish machine shop to construct the phosphorus atom models needed for short sections of the sugar-phosphate backbone found in DNA. Once these became available, we tested different ways the backbones might twist around each other in the center of the DNA molecule. Their regular repeating atomic structure should allow the atoms to come together in a consistent, repeated conformation. Following Wilkins's hunch, we focused on three-chain models. When one of these appeared to be almost plausible, Crick made a phone call to Wilkins to announce we had a model we thought might be DNA.

\smallskip
\noindent\textbf{中文。} 我们一回到剑桥，我就安排卡文迪许机修车间制作DNA中糖-磷酸骨架短段所需的磷原子模型。模型做好后，我们测试骨架在DNA分子中心彼此缠绕的不同方式。它们规则重复的原子结构应当使原子以一致、重复的构象聚在一起。按照威尔金斯的直觉，我们集中研究三链模型。当其中一个模型看上去几乎可行时，克里克给威尔金斯打电话，宣布我们有了一个自认为可能是DNA的模型。
\end{quote}

\begin{quote}
\noindent\textbf{English.} The next day both Wilkins and Franklin came up to see what we had done. The threat of unanticipated competition briefly united them in common purpose. Franklin wasted no time in faulting our basic concept. My memory was that she had reported almost no water present in crystalline DNA. In fact, the opposite was true. Being a crystallographic novice, I had confused the terms "unit cell" and "asymmetric unit." Crystalline DNA was in fact water-rich. Consequently, Franklin pointed out, the backbone had to be on the outside and not, as we had it, in the center, if only to accommodate all the water molecules she had observed in her crystals.

\smallskip
\noindent\textbf{中文。} 第二天，威尔金斯和富兰克林都来查看我们做了什么。意外竞争的威胁，短暂地把他们团结在共同目标上。富兰克林立刻指出我们基本概念的错误。我记得她报告过晶态DNA几乎不含水。事实上正相反。作为晶体学新手，我把“晶胞”和“不对称单元”混淆了。晶态DNA实际上富含水分。因此，富兰克林指出，骨架必须在外侧，而不能像我们的模型那样在中心，哪怕只是为了容纳她在晶体中观察到的所有水分子。
\end{quote}

\begin{quote}
\noindent\textbf{English.} That unfortunate November day cast a very long shadow. Franklin's opposition to model-building was reinforced. Doing experiments, not playing with Tinkertoy representations of atoms, was the way she intended to proceed. Even worse, Sir Lawrence Bragg passed down the word that Crick and I should desist from all further attempts at building a DNA model. It was further decreed that DNA research should be left to the King's lab, with Cambridge continuing to focus solely on proteins. There was no sense in two MRC-funded labs competing against each other. With no more bright ideas up our sleeves, Crick and I were reluctantly forced to back off, at least for the time being.

\smallskip
\noindent\textbf{中文。} 那个不幸的十一月日子投下了很长的阴影。富兰克林反对建模的立场得到强化。她打算继续的方式，是做实验，而不是玩用儿童拼插玩具代表原子的游戏。更糟的是，劳伦斯·布拉格爵士传话下来，要求克里克和我停止一切进一步构建DNA模型的尝试。还进一步规定，DNA研究应交给国王学院实验室，剑桥继续只专注蛋白质。两个由MRC资助的实验室相互竞争没有意义。我们袖子里也没有更多妙计，克里克和我只好不情愿地退后一步，至少暂时如此。
\end{quote}

\begin{quote}
\noindent\textbf{English.} It was not a good moment to be condemned to the DNA sidelines. Linus Pauling had written Wilkins to request a copy of the crystalline DNA diffraction pattern. Though Wilkins had declined, saying he wanted more time to interpret it himself, Pauling was hardly obliged to depend upon data from King's. If he wished, he could easily start serious X-ray diffraction studies at Caltech.

\smallskip
\noindent\textbf{中文。} 被判到DNA研究边线旁，这不是一个好时机。莱纳斯·鲍林已经写信给威尔金斯，请求一份晶态DNA衍射图样。威尔金斯拒绝了，说自己还需要更多时间解释它，但鲍林几乎不必依赖国王学院的数据。只要愿意，他完全可以在加州理工轻易开展严肃的X射线衍射研究。
\end{quote}

\begin{quote}
\noindent\textbf{English.} The following spring, I duly turned away from DNA and set about extending prewar studies on the pencil-shaped tobacco mosaic virus using the Cavendish's powerful new X-ray beam. This light experimental workload gave me plenty of time to wander through various Cambridge libraries. In the zoology building, I read Erwin Chargaff's paper describing his finding that the DNA bases adenine and thymine occurred in roughly equal amounts, as did the bases guanine and cytosine. Hearing of these one-to-one ratios Crick wondered whether, during DNA duplication, adenine residues might be attracted to thymine and vice versa, and whether a corresponding attraction might exist between guanine and cytosine. If so, base sequences on the "parental" chains, for example ATGC, would have to be complementary to those on "daughter" strands, yielding in this case TACG.

\smallskip
\noindent\textbf{中文。} 第二年春天，我如命令所要求的那样转离DNA，开始使用卡文迪许强大的新X射线束，扩展战前对铅笔状烟草花叶病毒的研究。这个实验工作量不重，使我有大量时间在剑桥各个图书馆游荡。在动物学楼里，我读到了埃尔温·查伽夫的论文，文中描述了他的发现：DNA碱基腺嘌呤和胸腺嘧啶的含量大致相等，鸟嘌呤和胞嘧啶的含量也大致相等。听到这些一比一比例后，克里克想知道，在DNA复制过程中，腺嘌呤残基是否可能吸引胸腺嘧啶，反之亦然；鸟嘌呤和胞嘧啶之间是否也可能存在相应吸引。如果是这样，“亲本”链上的碱基序列，例如ATGC，就必须与“子代”链上的序列互补，在这个例子中产生TACG。
\end{quote}

\begin{quote}
\noindent\textbf{English.} These remained idle thoughts until Erwin Chargaff came through Cambridge in the summer of 1952 on his way to the International Biochemical Congress in Paris. Chargaff expressed annoyance that neither Crick nor I saw the need to know the chemical structures of the four bases. He was even more upset when we told him that we could simply look up the structures in textbooks as the need arose. I was left hoping that Chargaff's data would prove irrelevant. Crick, however, was energized to do several experiments looking for molecular "sandwiches" that might form when adenine and thymine, or alternatively guanine and cytosine, were mixed together in solution. But his experiments went nowhere.

\smallskip
\noindent\textbf{中文。} 这些想法一直只是闲想，直到1952年夏天埃尔温·查伽夫途经剑桥，前往巴黎参加国际生物化学大会。查伽夫对克里克和我都不认为有必要知道四种碱基的化学结构感到恼火。当我们告诉他，若有需要可以直接去教科书上查结构时，他更加不悦。我只希望查伽夫的数据最终会被证明无关紧要。克里克却因此振奋起来，做了几项实验，寻找腺嘌呤和胸腺嘧啶，或者鸟嘌呤和胞嘧啶，在溶液中混合时可能形成的分子“三明治”。但他的实验没有进展。
\end{quote}

\begin{quote}
\noindent\textbf{English.} Like Chargaff, Linus Pauling also attended the International Biochemical Congress, where the big news was the latest result from the Phage Group. Alfred Hershey and Martha Chase at Cold Spring Harbor had just confirmed Avery's transforming principle: DNA was the hereditary material! Hershey and Chase proved that only the DNA of the phage virus enters bacterial cells; its protein coat remains on the outside. It was more obvious than ever that DNA must be understood at the molecular level if we were to uncover the essence of the gene. With Hershey and Chase's result the talk of the town, I was sure that Pauling would now bring his formidable intellect and chemical wisdom to bear on the problem of DNA.

\smallskip
\noindent\textbf{中文。} 和查伽夫一样，莱纳斯·鲍林也参加了国际生物化学大会。大会上的大新闻是噬菌体小组的最新结果。冷泉港的阿尔弗雷德·赫尔希和玛莎·蔡斯刚刚证实了艾弗里的转化原则：DNA就是遗传物质！赫尔希和蔡斯证明，只有噬菌体病毒的DNA进入细菌细胞；其蛋白质外壳留在外面。若要揭示基因的本质，就必须在分子层面理解DNA，这比以往任何时候都更明显。赫尔希和蔡斯的结果成为全城热议，我确信鲍林如今会把他强大的智力和化学智慧用于DNA问题。
\end{quote}

\begin{quote}
\noindent\textbf{English.} Early in 1953, Pauling did indeed publish a paper outlining the structure of DNA. Reading it anxiously I saw that he was proposing a three-chain model with sugar-phosphate backbones forming a dense central core. Superficially it was similar to our botched model of fifteen months earlier. But instead of using positively charged atoms, for example \(Mg^{2+}\), to stabilize the negatively charged backbones, Pauling made the unorthodox suggestion that the phosphates were held together by hydrogen bonds. But it seemed to me, the biologist, that such hydrogen bonds required extremely acidic conditions never found in cells. With a mad dash to Alexander Todd's nearby organic chemistry lab my belief was confirmed: the impossible had happened. The world's best-known, if not best, chemist had gotten his chemistry wrong. In effect, Pauling had knocked the A off of DNA. Our quarry was deoxyribonucleic acid, but the structure he was proposing was not even acidic.

\smallskip
\noindent\textbf{中文。} 1953年初，鲍林确实发表了一篇概述DNA结构的论文。我焦急地读着，看到他提出的是三链模型，糖-磷酸骨架形成致密中央核心。从表面看，这与我们十五个月前搞砸的模型相似。但鲍林没有使用带正电的原子，例如 \(Mg^{2+}\)，来稳定带负电的骨架，而是提出一个非正统建议：磷酸基团由氢键维系在一起。可在我这个生物学家看来，这样的氢键需要细胞内从不存在的极端酸性条件。我狂奔到附近亚历山大·托德的有机化学实验室，我的判断得到证实：不可能发生的事发生了。世界上最有名，即使不是最好的，化学家把化学弄错了。实际上，鲍林把DNA中的“A”去掉了。我们的猎物是脱氧核糖核酸，但他提出的结构甚至不是酸性的。
\end{quote}

\begin{quote}
\noindent\textbf{English.} Hurriedly I took the manuscript to London to inform Wilkins and Franklin they were still in the game. Convinced that DNA was not a helix, Franklin had no wish even to read the article and deal with the distraction of Pauling's helical ideas, even when I offered Crick's arguments for helices. Wilkins, however, was very interested indeed in the news I brought; he was now more certain than ever that DNA was helical. To prove the point, he showed me a photograph obtained more than six months earlier by Franklin's graduate student Raymond Gosling, who had X-rayed the so-called B form of DNA. Until that moment, I did not know a B form even existed. Franklin had put this picture aside, preferring to concentrate on the A form, which she thought would more likely yield useful data. The X-ray pattern of this B form was a distinct cross. Since Crick and others had already deduced that such a pattern of reflections would be created by a helix, this evidence made it clear that DNA had to be a helix! In fact, despite Franklin's reservations, this was no surprise. Geometry itself suggested that a helix was the most logical arrangement for a long string of repeating units such as the nucleotides of DNA. But we still did not know what that helix looked like, nor how many chains it contained.

\smallskip
\noindent\textbf{中文。} 我匆忙把手稿带到伦敦，告诉威尔金斯和富兰克林他们还在游戏中。富兰克林确信DNA不是螺旋，甚至不愿读这篇文章，不愿应付鲍林螺旋思想的干扰，即使我提出克里克支持螺旋的论据也无济于事。然而，威尔金斯对我带来的消息确实非常感兴趣；他现在比以往更确信DNA是螺旋。为了证明这一点，他给我看了一张六个多月前由富兰克林研究生雷蒙德·戈斯林获得的照片，戈斯林曾对所谓DNA B型进行X射线照射。在那一刻之前，我甚至不知道B型存在。富兰克林把这张照片搁在一旁，更愿意集中研究A型，因为她认为A型更可能产生有用数据。这个B型的X射线图案呈现清晰十字。克里克和其他人已经推导出，螺旋会产生这种反射图案，因此这一证据清楚表明DNA必定是螺旋！事实上，尽管富兰克林有所保留，这并不令人意外。几何学本身提示，对DNA核苷酸这样长串重复单元来说，螺旋是最合乎逻辑的排列。但我们仍不知道这个螺旋是什么样子，也不知道它包含几条链。
\end{quote}

\begin{quote}
\noindent\textbf{English.} The time had come to resume building helical models of DNA. Pauling was bound to realize soon enough that his brainchild was wrong. I urged Wilkins to waste no time. But he wanted to wait until Franklin had completed her scheduled departure for another lab later that spring. She had decided to move on to avoid the unpleasantness at King's. Before leaving, she had been ordered to stop further work with DNA and had already passed on many of her diffraction images to Wilkins.

\smallskip
\noindent\textbf{中文。} 重新构建DNA螺旋模型的时刻到了。鲍林迟早会意识到自己的心血结晶是错的。我敦促威尔金斯不要浪费时间。但他想等到富兰克林按计划在那年晚春离开，前往另一间实验室之后。她已经决定离开，以避开国王学院的不愉快。离开前，她被命令停止进一步研究DNA，并已把许多衍射图像交给威尔金斯。
\end{quote}

\begin{figure}[htbp]
\centering
\ocrimage[width=0.45\linewidth]{images/19c9e38243b212a63b0c2c3bdaff52a3070510c4e38cefa3383690dda5682970.jpg}
\hfill
\ocrimage[width=0.45\linewidth]{images/e490980fa1fb6465bd4c9013ed309954bdf81cfd15b9d02dcee06b0de3b45061.jpg}
\caption{X-ray photos of the A and B forms of DNA from, respectively, Maurice Wilkins and Rosalind Franklin. The differences in molecular structure are caused by differences in the amount of water associated with each DNA molecule.\\DNA的A型和B型X射线照片，分别来自莫里斯·威尔金斯和罗莎琳德·富兰克林。分子结构差异由每个DNA分子结合水量的差异造成。}
\end{figure}

\begin{quote}
\noindent\textbf{English.} When I returned to Cambridge and broke the news of the DNA B form, Bragg no longer saw any reason for Crick and me to avoid DNA. He very much wanted the DNA structure to be found on his side of the Atlantic. So we went back to model-building, looking for a way the known basic components of DNA---the backbone of the molecule and the four different bases, adenine, thymine, guanine, and cytosine---could fit together to make a helix. I commissioned the shop at the Cavendish to make us a set of tin bases, but they could not produce them fast enough for me: I ended up cutting out rough approximations from stiff cardboard.

\smallskip
\noindent\textbf{中文。} 我回到剑桥并透露DNA B型的消息后，布拉格再也看不到克里克和我避开DNA的理由。他非常希望DNA结构能在大西洋这一侧被发现。于是我们重新开始建模，寻找一种方式，让DNA已知的基本组成部分，即分子骨架和四种不同碱基腺嘌呤、胸腺嘧啶、鸟嘌呤、胞嘧啶，能够组合成螺旋。我委托卡文迪许车间为我们制作一套锡制碱基，但他们做得不够快；最后我用硬纸板剪出粗略近似物。
\end{quote}

\begin{quote}
\noindent\textbf{English.} By this time I realized the DNA density-measurement evidence actually slightly favored a two-chain, rather than three-chain, model. So I decided to search out plausible double helices. As a biologist, I preferred the idea of a genetic molecule made of two, rather than three, components. After all, chromosomes, like cells, increase in number by duplicating, not triplicating.

\smallskip
\noindent\textbf{中文。} 到这时，我意识到DNA密度测量证据实际上稍微支持两链模型，而不是三链模型。因此，我决定寻找可行的双螺旋。作为生物学家，我更喜欢遗传分子由两个组成部分而非三个组成部分构成的想法。毕竟，染色体像细胞一样，是通过复制而不是三倍化来增加数目的。
\end{quote}

\begin{figure}[htbp]
\centering
\ocrimage[width=0.70\linewidth]{images/c879c31c90a3b1b3c2b93ae578d9b3f6990843106b896ca71172ac4499e54282.jpg}
\caption{The chemical backbone of DNA.\\DNA的化学骨架。}
\end{figure}

\begin{quote}
\noindent\textbf{English.} I knew that our previous model with the backbone on the inside and the bases hanging out was wrong. Chemical evidence from the University of Nottingham, which I had too long ignored, indicated that the bases must be hydrogen-bonded to each other. They could only form bonds like this in the regular manner implied by the X-ray diffraction data if they were in the center of the molecule. But how could they come together in pairs? For two weeks I got nowhere, misled by an error in my nucleic acid chemistry textbook. Happily, on February 27, Jerry Donohue, a theoretical chemist visiting the Cavendish from Caltech, pointed out that the textbook was wrong. So I changed the locations of the hydrogen atoms on my cardboard cutouts of the molecules.

\smallskip
\noindent\textbf{中文。} 我知道我们先前那个骨架在内、碱基朝外的模型是错的。我长期忽视的诺丁汉大学化学证据表明，碱基之间必须彼此形成氢键。只有当碱基位于分子中心时，它们才能以X射线衍射数据所暗示的规则方式形成这种键。可是它们如何成对聚合？两个星期里我毫无进展，被核酸化学教科书中的一个错误误导了。幸运的是，2月27日，从加州理工来卡文迪许访问的理论化学家杰里·多诺休指出教科书错了。于是我改变了分子硬纸板剪片上氢原子的位置。
\end{quote}

\begin{quote}
\noindent\textbf{English.} The next morning, February 28, 1953, the key features of the DNA model all fell into place. The two chains were held together by strong hydrogen bonds between adenine--thymine and guanine--cytosine base pairs. The inferences Crick had drawn the year before based on Chargaff's research had indeed been correct. Adenine does bond to thymine and guanine does bond to cytosine, but not through flat surfaces to form molecular sandwiches. When Crick arrived, he took it all in rapidly, and gave my base-pairing scheme his blessing. He realized right away that it would result in the two strands of the double helix running in opposite directions.

\smallskip
\noindent\textbf{中文。} 第二天早晨，1953年2月28日，DNA模型的关键特征全部就位。两条链由腺嘌呤--胸腺嘧啶和鸟嘌呤--胞嘧啶碱基对之间的强氢键维系。克里克前一年基于查伽夫研究作出的推论确实正确。腺嘌呤会与胸腺嘧啶结合，鸟嘌呤会与胞嘧啶结合，但不是通过平面表面形成分子三明治。克里克到来后，很快理解了一切，并认可了我的碱基配对方案。他立刻意识到，这将导致双螺旋的两条链以相反方向运行。
\end{quote}

\begin{quote}
\noindent\textbf{English.} It was quite a moment. We felt sure that this was it. Anything that simple, that elegant just had to be right. What got us most excited was the complementarity of the base sequences along the two chains. If you knew the sequence---the order of bases---along one chain, you automatically knew the sequence along the other. It was immediately apparent that this must be how the genetic messages of genes are copied so exactly when chromosomes duplicate prior to cell division. The molecule would "unzip" to form two separate strands. Each separate strand then could serve as the template for the synthesis of a new strand, one double helix becoming two.

\smallskip
\noindent\textbf{中文。} 那真是一个了不起的时刻。我们确信就是它了。如此简单、如此优雅的东西，必然是正确的。最令我们兴奋的是两条链上碱基序列的互补性。如果你知道一条链上的序列，也就是碱基顺序，就会自动知道另一条链上的序列。立刻显而易见的是，这一定就是染色体在细胞分裂前复制时，基因的遗传信息能够如此精确复制的方式。这个分子会像拉链一样“拉开”，形成两条单独链。每条单链随后都可以作为模板合成一条新链，一个双螺旋变成两个。
\end{quote}

\begin{quote}
\noindent\textbf{English.} In \emph{What Is Life?} Schrödinger had suggested that the language of life might be like Morse code, a series of dots and dashes. He was not far off. The language of DNA is a linear series of As, Ts, Gs, and Cs. And just as transcribing a page out of a book can result in the odd typo, the rare mistake creeps in when all these As, Ts, Gs, and Cs are being copied along a chromosome. These errors are the mutations geneticists had talked about for almost fifty years. Change an "i" to an "a" and "Jim" becomes "Jam" in English; change a T to a C and "ATG" becomes "ACG" in DNA.

\smallskip
\noindent\textbf{中文。} 在《生命是什么？》中，薛定谔曾提出，生命的语言也许像莫尔斯电码，是一系列点和划。他离真相并不远。DNA的语言是由A、T、G、C组成的线性序列。正如抄写书页会偶尔出现错别字，当所有这些A、T、G、C沿染色体被复制时，罕见错误也会潜入。这些错误就是遗传学家谈论了近五十年的突变。在英语中，把“Jim”里的“i”改成“a”，“Jim”就变成“Jam”；在DNA中，把T改成C，“ATG”就变成“ACG”。
\end{quote}

\begin{figure}[htbp]
\centering
\ocrimage[width=0.45\linewidth]{images/9ed2d5675b2303884242823021aa3e3ba5819dec62e874dbeaa1906ab110f706.jpg}
\hfill
\ocrimage[width=0.45\linewidth]{images/c99acaa9355ddb2b918c14b2bb9ebb29ed0780f6348debfe1bfe251cb90b9209.jpg}
\caption{The insight that made it all come together: complementary pairing of the bases.\\让一切贯通的洞见：碱基的互补配对。}
\end{figure}

\begin{figure}[htbp]
\centering
\ocrimage[width=0.45\linewidth]{images/89cf17fe3c783f0129e6decdcfd90773809f800d82fac3712a7aa030f95c657c.jpg}
\hfill
\ocrimage[width=0.45\linewidth]{images/556ed40ecd12f173f0a6361a38347c46ec70df5becf9e667c1cb7ac95a6ed493.jpg}
\caption{Bases and backbone in place: the double helix. (A) is a schematic showing the system of base-pairing that binds the two strands together. (B) is a "space-filling" model showing, to scale, the atomic detail of the molecule.\\碱基和骨架各就各位：双螺旋。（A）为示意图，显示把两条链结合在一起的碱基配对系统。（B）为“空间填充”模型，按比例展示分子的原子细节。}
\end{figure}

\begin{quote}
\noindent\textbf{English.} The double helix made sense chemically and it made sense biologically. Now there was no need to be concerned about Schrödinger's suggestion that new laws of physics might be necessary for an understanding of how the hereditary code-script is duplicated: genes in fact were no different from the rest of chemistry. Later that day, during lunch at the Eagle, the pub virtually adjacent to the Cavendish Lab, Crick, ever the talker, could not help but tell everyone we had just found the "secret of life." I myself, though no less electrified by the thought, would have waited until we had a pretty three-dimensional model to show off.

\smallskip
\noindent\textbf{中文。} 双螺旋在化学上说得通，在生物学上也说得通。现在无需再担心薛定谔关于理解遗传密码脚本复制也许需要新物理定律的说法：基因事实上并不不同于其他化学。当天晚些时候，在几乎紧挨卡文迪许实验室的“鹰”酒馆吃午饭时，一贯健谈的克里克忍不住告诉所有人，我们刚刚发现了“生命的秘密”。至于我自己，虽然同样为这个想法激动不已，却会等到我们有了漂亮的三维模型可展示之后再说。
\end{quote}

\begin{quote}
\noindent\textbf{English.} Among the first to see our demonstration model was the chemist Alexander Todd. That the nature of the gene was so simple both surprised and pleased him.

\smallskip
\noindent\textbf{中文。} 最早看到我们演示模型的人之一，是化学家亚历山大·托德。基因的本质如此简单，既让他惊讶，也让他高兴。
\end{quote}

\begin{quote}
\noindent\textbf{English.} Later, however, he must have asked himself why his own lab, having established the general chemical structure of DNA chains, had not moved on to asking how the chains folded up in three dimensions. Instead the essence of the molecule was left to be discovered by a two-man team, a biologist and a physicist, neither of whom possessed a detailed command even of undergraduate chemistry. But paradoxically, this was, at least in part, the key to our success: Crick and I arrived at the double helix first precisely because most chemists at that time thought DNA too big a molecule to understand by chemical analysis.

\smallskip
\noindent\textbf{中文。} 然而后来，他一定问过自己：既然自己的实验室已经建立了DNA链的一般化学结构，为什么没有进一步追问这些链在三维空间中如何折叠？相反，这个分子的本质留给了两人小组来发现，一个生物学家和一个物理学家，而两人甚至都没有详细掌握本科化学。但矛盾的是，这至少部分正是我们成功的关键：克里克和我之所以最先抵达双螺旋，恰恰因为当时多数化学家认为DNA分子太大，无法用化学分析来理解。
\end{quote}

\begin{quote}
\noindent\textbf{English.} At the same time, the only two chemists with the vision to seek DNA's 3-D structure made major tactical mistakes: Rosalind Franklin's was her resistance to model-building; Linus Pauling's was a matter of simply neglecting to read the existing literature on DNA, particularly the data on its base composition published by Chargaff. Ironically, Pauling and Chargaff sailed across the Atlantic on the same ship following the Paris Biochemical Congress in 1952, but failed to hit it off. Pauling was long accustomed to being right. And he believed there was no chemical problem he could not work out from first principles by himself. Usually this confidence was not misplaced. During the Cold War, as a prominent critic of the American nuclear weapons development program, he was questioned by the FBI after giving a talk. How did he know how much plutonium there is in an atomic bomb? Pauling's response was "Nobody told me. I figured it out."

\smallskip
\noindent\textbf{中文。} 与此同时，唯有两个有远见去寻找DNA三维结构的化学家，都犯下了重大战术错误：罗莎琳德·富兰克林的错误在于抵制模型构建；莱纳斯·鲍林的错误则只是没有阅读已有DNA文献，尤其是查伽夫发表的碱基组成数据。具有讽刺意味的是，1952年巴黎生物化学大会后，鲍林和查伽夫乘同一艘船横渡大西洋，却没能投缘。鲍林早已习惯于自己是对的。他相信没有任何化学问题不能由自己从第一原理推出来。通常，这种自信并没有错置。冷战期间，他作为美国核武器发展计划的著名批评者，在一次演讲后受到FBI询问。他怎么知道一枚原子弹里有多少钚？鲍林回答说：“没人告诉我。我算出来的。”
\end{quote}

\begin{quote}
\noindent\textbf{English.} Over the next several months Crick and, to a lesser extent, I relished showing off our model to an endless stream of curious scientists. However, the Cambridge biochemists did not invite us to give a formal talk in the biochemistry building. They started to refer to it as the "WC," punning our initials with those used in Britain for the toilet or water closet. That we had found the double helix without doing experiments irked them.

\smallskip
\noindent\textbf{中文。} 在接下来的几个月里，克里克和我，后者程度稍弱，都非常享受向源源不断的好奇科学家展示我们的模型。然而，剑桥的生物化学家并没有邀请我们到生化楼作正式报告。他们开始把它称作“WC”，用我们的姓名首字母与英国表示厕所或盥洗室的缩写双关。我们没有做实验却找到了双螺旋，这让他们恼火。
\end{quote}

\begin{quote}
\noindent\textbf{English.} The manuscript that we submitted to \emph{Nature} in early April was published just over three weeks later, on April 25, 1953. Accompanying it were two longer papers by Franklin and Wilkins, both supporting the general correctness of our model. In June, I gave the first presentation of our model at the Cold Spring Harbor symposium on viruses. Max Delbrück saw to it that I was offered, at the last minute, an invitation to speak. To this intellectually high-powered meeting I brought a three-dimensional model built in the Cavendish, the adenine--thymine base pairs in red and the guanine--cytosine base pairs in green.

\smallskip
\noindent\textbf{中文。} 我们四月初提交给《自然》的手稿，在三周多后，即1953年4月25日发表。与之同期刊出的还有富兰克林和威尔金斯的两篇更长论文，二者都支持我们模型的大体正确性。六月，我在冷泉港病毒专题讨论会上首次介绍了我们的模型。马克斯·德尔布吕克确保我在最后时刻获得了发言邀请。面对这场智力强度极高的会议，我带来了一个在卡文迪许制作的三维模型，其中腺嘌呤--胸腺嘧啶碱基对为红色，鸟嘌呤--胞嘧啶碱基对为绿色。
\end{quote}

\section*{Molecular Structure of Nucleic Acids}
\begin{center}
\large A Structure for Deoxyribose Nucleic Acid\\
\large 核酸的分子结构：脱氧核糖核酸的一种结构
\end{center}

\begin{quote}
\noindent\textbf{English.} We wish to suggest a structure for the salt of deoxyribose nucleic acid (D.N.A.). This structure has novel features which are of considerable biological interest.

\smallskip
\noindent\textbf{中文。} 我们希望提出脱氧核糖核酸（D.N.A.）盐的一种结构。这个结构具有若干新颖特征，具有相当大的生物学意义。
\end{quote}

\begin{quote}
\noindent\textbf{English.} A structure for nucleic acid has already been proposed by Pauling and Corey\textsuperscript{1}. They kindly made their manuscript available to us in advance of publication. Their model consists of three intertwined chains, with the phosphates near the fibre axis, and the bases on the outside. In our opinion, this structure is unsatisfactory for two reasons: (1) We believe that the material which gives the X-ray diagrams is the salt, not the free acid. Without the acidic hydrogen atoms it is not clear what forces would hold the structure together, especially as the negatively charged phosphates near the axis will repel each other. (2) Some of the van der Waals distances appear to be too small.

\smallskip
\noindent\textbf{中文。} 鲍林和科里已经提出过一种核酸结构\textsuperscript{1}。他们友好地在发表前把手稿提供给我们。他们的模型由三条相互缠绕的链组成，磷酸基靠近纤维轴，碱基在外侧。我们认为这个结构有两个方面不能令人满意：（1）我们相信产生X射线图样的材料是盐，而不是游离酸。没有酸性氢原子，就不清楚是什么力量把结构维系在一起，尤其是靠近轴心的带负电磷酸基会彼此排斥。（2）某些范德华距离似乎太小。
\end{quote}

\begin{quote}
\noindent\textbf{English.} Another three-chain structure has also been suggested by Fraser, in the press. In his model the phosphates are on the outside and the bases on the inside, linked together by hydrogen bonds. This structure as described is rather ill-defined, and for this reason we shall not comment on it.

\smallskip
\noindent\textbf{中文。} 弗雷泽也提出了另一种三链结构，论文即将发表。在他的模型中，磷酸基在外侧，碱基在内侧，并由氢键连接在一起。按其描述，这一结构相当不明确，因此我们不作评论。
\end{quote}

\begin{figure}[htbp]
\centering
\ocrimage[width=0.55\linewidth]{images/a43c7f1bb5bbda79e9e814375d12af675f324f9ef0f2b73d1f00ab1a62b64b57.jpg}
\caption{This figure is purely diagrammatic. The two ribbons symbolize the two phosphate--sugar chains, and the horizontal rods the pairs of bases holding the chains together. The vertical line marks the fibre axis.\\此图纯属示意。两条带状物表示两条磷酸--糖链，水平短杆表示把链维系在一起的碱基对。垂直线标示纤维轴。}
\end{figure}

\begin{quote}
\noindent\textbf{English.} We wish to put forward a radically different structure for the salt of deoxyribose nucleic acid. This structure has two helical chains each coiled round the same axis (see diagram). We have made the usual chemical assumptions, namely, that each chain consists of phosphate diester groups joining \(\beta\)-D-deoxyribofuranose residues with 3',5' linkages. The two chains, but not their bases, are related by a dyad perpendicular to the fibre axis. Both chains follow right-handed helices, but owing to the dyad the sequences of the atoms in the two chains run in opposite directions. Each chain loosely resembles Furberg's model No. 1; that is, the bases are on the inside of the helix and the phosphates on the outside. The configuration of the sugar and the atoms near it is close to Furberg's "standard configuration," the sugar being roughly perpendicular to the attached base. There is a residue on each chain every 3.4 \AA{} in the \(z\)-direction. We have assumed an angle of \(36^\circ\) between adjacent residues in the same chain, so that the structure repeats after 10 residues on each chain, that is, after 34 \AA{}. The distance of a phosphorus atom from the fibre axis is 10 \AA{}. As the phosphates are on the outside, cations have easy access to them.

\smallskip
\noindent\textbf{中文。} 我们希望为脱氧核糖核酸盐提出一种根本不同的结构。这个结构有两条螺旋链，各自围绕同一轴缠绕（见图）。我们采用通常的化学假设，即每条链由磷酸二酯基团组成，这些基团通过3',5'键连接 \(\beta\)-D-脱氧核糖呋喃糖残基。两条链（但不包括它们的碱基）由一条垂直于纤维轴的二重轴相关。两条链都遵循右手螺旋，但由于二重轴，两条链中原子的序列方向相反。每条链大致类似富伯格的1号模型；也就是说，碱基在螺旋内部，磷酸基在外部。糖及其附近原子的构型接近富伯格的“标准构型”，糖大致垂直于所连接的碱基。沿 \(z\) 方向，每条链每隔3.4埃有一个残基。我们假定同一链上相邻残基之间的角度为 \(36^\circ\)，因此每条链上10个残基后结构重复，即34埃后重复。磷原子距纤维轴10埃。由于磷酸基在外侧，阳离子很容易接近它们。
\end{quote}

\begin{quote}
\noindent\textbf{English.} The structure is an open one, and its water content is rather high. At lower water contents we would expect the bases to tilt so that the structure could become more compact.

\smallskip
\noindent\textbf{中文。} 该结构是开放的，含水量相当高。在含水量较低时，我们预期碱基会倾斜，使结构变得更加紧凑。
\end{quote}

\begin{quote}
\noindent\textbf{English.} The novel feature of the structure is the manner in which the two chains are held together by the purine and pyrimidine bases. The planes of the bases are perpendicular to the fibre axis. They are joined together in pairs, a single base from one chain being hydrogen-bonded to a single base from the other chain, so that the two lie side by side with identical \(z\)-coordinates. One of the pair must be a purine and the other a pyrimidine for bonding to occur. The hydrogen bonds are made as follows: purine position 1 to pyrimidine position 1; purine position 6 to pyrimidine position 6.

\smallskip
\noindent\textbf{中文。} 这一结构的新颖特征，在于两条链由嘌呤和嘧啶碱基维系在一起的方式。碱基平面垂直于纤维轴。它们成对连接，一条链上的一个碱基与另一条链上的一个碱基形成氢键，使二者并排处在相同的 \(z\) 坐标上。要发生成键，一对中一个必须是嘌呤，另一个必须是嘧啶。氢键形成方式如下：嘌呤1位连到嘧啶1位；嘌呤6位连到嘧啶6位。
\end{quote}

\begin{quote}
\noindent\textbf{English.} If it is assumed that the bases only occur in the structure in the most plausible tautomeric forms, that is, with the keto rather than the enol configurations, it is found that only specific pairs of bases can bond together. These pairs are: adenine, a purine, with thymine, a pyrimidine; and guanine, a purine, with cytosine, a pyrimidine.

\smallskip
\noindent\textbf{中文。} 如果假定碱基在结构中只以最合理的互变异构形式出现，也就是以酮式而不是烯醇式构型出现，就会发现只有特定碱基对能够彼此成键。这些配对是：腺嘌呤（嘌呤）与胸腺嘧啶（嘧啶）；鸟嘌呤（嘌呤）与胞嘧啶（嘧啶）。
\end{quote}

\begin{quote}
\noindent\textbf{English.} In other words, if an adenine forms one member of a pair, on either chain, then on these assumptions the other member must be thymine; similarly for guanine and cytosine. The sequence of bases on a single chain does not appear to be restricted in any way. However, if only specific pairs of bases can be formed, it follows that if the sequence of bases on one chain is given, then the sequence on the other chain is automatically determined.

\smallskip
\noindent\textbf{中文。} 换言之，如果腺嘌呤在任一条链上构成一对中的一个成员，那么在这些假设下，另一成员必定是胸腺嘧啶；鸟嘌呤与胞嘧啶同理。单条链上的碱基序列似乎不受任何限制。然而，如果只有特定碱基对能够形成，那么只要给定一条链上的碱基序列，另一条链上的序列就会自动确定。
\end{quote}

\begin{quote}
\noindent\textbf{English.} It has been found experimentally\textsuperscript{2,4} that the ratio of the amounts of adenine to thymine, and the ratio of guanine to cytosine, are always very close to unity for deoxyribose nucleic acid.

\smallskip
\noindent\textbf{中文。} 实验已经发现\textsuperscript{2,4}，在脱氧核糖核酸中，腺嘌呤与胸腺嘧啶的含量比，以及鸟嘌呤与胞嘧啶的含量比，总是非常接近1。
\end{quote}

\begin{quote}
\noindent\textbf{English.} It is probably impossible to build this structure with a ribose sugar in place of the deoxyribose, as the extra oxygen atom would make too close a van der Waals contact.

\smallskip
\noindent\textbf{中文。} 如果用核糖替代脱氧核糖，可能无法构建这一结构，因为额外的氧原子会造成过近的范德华接触。
\end{quote}

\begin{quote}
\noindent\textbf{English.} The previously published X-ray data\textsuperscript{5,6} on deoxyribose nucleic acid are insufficient for a rigorous test of our structure. So far as we can tell, it is roughly compatible with the experimental data, but it must be regarded as unproved until it has been checked against more exact results. Some of these are given in the following communications. We were not aware of the details of the results presented there when we devised our structure, which rests mainly though not entirely on published experimental data and stereochemical arguments.

\smallskip
\noindent\textbf{中文。} 先前发表的关于脱氧核糖核酸的X射线数据\textsuperscript{5,6}，不足以对我们的结构作严格检验。就我们所能判断，它大体上与实验数据相容，但在与更精确结果核对之前，必须视为尚未证明。随后的论文给出了一些这样的结果。我们设计该结构时，并不知道那些结果的细节；我们的结构主要但并非完全建立在已发表实验数据和立体化学论证之上。
\end{quote}

\begin{quote}
\noindent\textbf{English.} It has not escaped our notice that the specific pairing we have postulated immediately suggests a possible copying mechanism for the genetic material.

\smallskip
\noindent\textbf{中文。} 我们注意到，我们所假定的特异性配对会立即提示遗传物质一种可能的复制机制。
\end{quote}

\begin{quote}
\noindent\textbf{English.} Full details of the structure, including the conditions assumed in building it, together with a set of co-ordinates for the atoms, will be published elsewhere.

\smallskip
\noindent\textbf{中文。} 该结构的完整细节，包括构建时假定的条件以及一组原子坐标，将另文发表。
\end{quote}

\begin{quote}
\noindent\textbf{English.} We are much indebted to Dr. Jerry Donohue for constant advice and criticism, especially on interatomic distances. We have also been stimulated by a knowledge of the general nature of the unpublished experimental results and ideas of Dr. M. H. F. Wilkins, Dr. R. E. Franklin and their co-workers at King's College, London. One of us (J. D. W.) has been aided by a fellowship from the National Foundation for Infantile Paralysis.

\smallskip
\noindent\textbf{中文。} 我们非常感谢杰里·多诺休博士持续不断的建议和批评，尤其是在原子间距离方面。我们还受到伦敦国王学院M. H. F. 威尔金斯博士、R. E. 富兰克林博士及其同事未发表实验结果和想法的一般性质的启发。我们中的一人（J. D. W.）获得了国家小儿麻痹症基金会奖学金资助。
\end{quote}

\begin{flushright}
J. D. Watson\\
F. H. C. Crick
\end{flushright}

\begin{quote}
\noindent Medical Research Council Unit for the Study of the Molecular Structure of Biological Systems, Cavendish Laboratory, Cambridge. April 2.

\smallskip
\noindent 医学研究委员会生物系统分子结构研究单元，剑桥卡文迪许实验室。4月2日。
\end{quote}

\begin{enumerate}
\item Pauling, L., and Corey, R. B., \emph{Nature}, 171, 346 (1953); \emph{Proc. U.S. Nat. Acad. Sci.}, 39, 84 (1953).
\item Furberg, S., \emph{Acta Chem. Scand.}, 6, 634 (1952).
\item Chargaff, E., for references see Zamenhof, S., Brawerman, G., and Chargaff, E., \emph{Biochim. et Biophys. Acta}, 9, 402 (1952).
\item Wyatt, G. R., \emph{J. Gen. Physiol.}, 36, 201 (1952).
\item Astbury, W. T., \emph{Symp. Soc. Exp. Biol.} 1, \emph{Nucleic Acid}, 66 (Camb. Univ. Press, 1947).
\item Wilkins, M. H. F., and Randall, J. T., \emph{Biochim. et Biophys. Acta}, 10, 192 (1953).
\end{enumerate}

\begin{quote}
\noindent\textbf{English.} Short and sweet: our \emph{Nature} paper announcing the discovery. The same issue also carried longer articles by Rosalind Franklin and Maurice Wilkins.

\smallskip
\noindent\textbf{中文。} 简短而漂亮：这就是我们宣布发现的《自然》论文。同一期还刊登了罗莎琳德·富兰克林和莫里斯·威尔金斯的较长论文。
\end{quote}

\begin{figure}[htbp]
\centering
\ocrimage[width=0.68\linewidth]{images/5f5c9a86c1556b15bc87624993ecaf2a5350fa6baa71c3325042a808c1935d14.jpg}
\caption{Unveiling the double helix: my lecture at Cold Spring Harbor Laboratory, June 1953.\\揭开双螺旋：1953年6月，我在冷泉港实验室作报告。}
\end{figure}

\begin{quote}
\noindent\textbf{English.} In the audience was Seymour Benzer, yet another ex-physicist who had heeded the clarion call of Schrödinger's book. He immediately understood what our breakthrough meant for his studies of mutations in viruses. He realized that he could now do for a short stretch of bacteriophage DNA what Morgan's boys had done forty years earlier for fruit fly chromosomes: he would map mutations---determine their order---along a gene, just as the fruit fly pioneers had mapped genes along a chromosome. Like Morgan, Benzer would have to depend on recombination to generate new genetic combinations, but, whereas Morgan had the advantage of a ready mechanism of recombination---the production of sex cells in a fruit fly---Benzer had to induce recombination by simultaneously infecting a single bacterial host cell with two different strains of bacteriophage, which differed by one or more mutations in the region of interest. Within the bacterial cell, recombination---the exchange of segments of molecules---would occasionally occur between the different viral DNA molecules, producing new permutations of mutations, so-called "recombinants." Within a single astonishingly productive year in his Purdue University lab, Benzer produced a map of a single bacteriophage gene, rII, showing how a series of mutations---all errors in the genetic script---were laid out linearly along the virus DNA. The language was simple and linear, just like a line of text on the written page.

\smallskip
\noindent\textbf{中文。} 听众中有西摩·本泽，他也是一位响应薛定谔之书号角召唤的前物理学家。他立刻理解我们的突破对自己研究病毒突变意味着什么。他意识到，现在他可以对一小段噬菌体DNA做四十年前“摩尔根的男孩们”对果蝇染色体所做的事：他将沿一个基因绘制突变图谱，确定它们的顺序，正如果蝇先驱者沿染色体绘制基因图谱一样。和摩尔根一样，本泽必须依赖重组来产生新的遗传组合；但摩尔根拥有现成的重组机制，即果蝇性细胞的产生，而本泽必须通过用两种不同噬菌体菌株同时感染同一个细菌宿主细胞来诱导重组，这两种菌株在所关注区域相差一个或多个突变。在细菌细胞内，不同病毒DNA分子之间偶尔会发生重组，即分子片段交换，产生突变的新排列，也就是所谓“重组体”。在普渡大学实验室里惊人高产的一年中，本泽绘制出单个噬菌体基因rII的图谱，显示一系列突变，即遗传脚本中的错误，如何沿病毒DNA线性排列。这种语言简单而线性，就像书页上一行文字。
\end{quote}

\begin{quote}
\noindent\textbf{English.} The response of the Hungarian physicist Leo Szilard to my Cold Spring Harbor talk on the double helix was less academic. His question was, "Can you patent it?" At one time Szilard's main source of income had been a patent that he held with Einstein, and he had later tried unsuccessfully to patent with Enrico Fermi the nuclear reactor they built at the University of Chicago in 1942. But then as now patents were given only for useful inventions and at the time no one could conceive of a practical use for DNA. Perhaps then, Szilard suggested, we should copyright it.

\smallskip
\noindent\textbf{中文。} 匈牙利物理学家利奥·西拉德对我在冷泉港关于双螺旋报告的反应就不那么学术了。他的问题是：“你们能给它申请专利吗？”西拉德曾有一段时间主要收入来自他与爱因斯坦共同持有的一项专利，后来他还曾试图与恩里科·费米一起为他们1942年在芝加哥大学建造的核反应堆申请专利，但没有成功。然而无论当时还是现在，专利只授予有用发明，而当时无人能想象DNA有什么实际用途。那么，西拉德建议，也许我们应该给它登记版权。
\end{quote}

\begin{figure}[htbp]
\centering
\ocrimage[width=0.70\linewidth]{images/a026b9883dbbda8eac990954a994fc9eaad9b41e2e1663ff203da9d998f6ce50.jpg}
\caption{DNA replication: the double helix is unzipped and each strand copied.\\DNA复制：双螺旋被拉开，每条链都被复制。}
\end{figure}

\begin{quote}
\noindent\textbf{English.} There remained, however, a single missing piece in the double helical jigsaw puzzle: our unzipping idea for DNA replication had yet to be experimentally verified. Max Delbrück, for example, was unconvinced. Though he liked the double helix as a model, he worried that unzipping it might generate horrible knots. Five years later, a former student of Pauling's, Matt Meselson, and the equally bright young phage worker Frank Stahl put to rest such fears when they published the results of a single elegant experiment.

\smallskip
\noindent\textbf{中文。} 然而，双螺旋拼图中仍缺少一块：我们关于DNA复制时“拉开拉链”的想法尚未得到实验证实。例如马克斯·德尔布吕克并不信服。虽然他喜欢双螺旋作为模型，却担心拉开它可能产生可怕的结。五年后，鲍林昔日学生马特·梅塞尔森和同样聪明的年轻噬菌体研究者弗兰克·斯塔尔，发表了一项优雅实验的结果，平息了这种担忧。
\end{quote}

\begin{quote}
\noindent\textbf{English.} They had met in the summer of 1954 at the Marine Biological Laboratory at Woods Hole, Massachusetts, where I was then lecturing, and agreed---over a good many gin martinis---that they should get together to do some science. The result of their collaboration has been described as "the most beautiful experiment in biology."

\smallskip
\noindent\textbf{中文。} 他们1954年夏天在马萨诸塞州伍兹霍尔海洋生物学实验室相遇，当时我在那里讲课。两人在许多杯金酒马提尼之后达成共识：他们应该合作做点科学。他们合作的结果后来被称为“生物学中最美丽的实验”。
\end{quote}

\begin{figure}[htbp]
\centering
\ocrimage[width=0.72\linewidth]{images/72946d8f3811ea6bdd57f651161bc59556e5a1d901a9dae3b0247b825b8790ac.jpg}
\caption{The Meselson--Stahl experiment.\\梅塞尔森--斯塔尔实验。}
\end{figure}

\begin{quote}
\noindent\textbf{English.} They used a centrifugation technique that allowed them to sort molecules according to slight differences in weight; following a centrifugal spin, heavier molecules end up nearer the bottom of the test tube than lighter ones. Because nitrogen atoms (N) are a component of DNA, and because they exist in two distinct forms, one light and one heavy, Meselson and Stahl were able to tag segments of DNA and thereby track the process of its replication in bacteria. Initially all the bacteria were raised in a medium containing heavy N, which was thus incorporated in both strands of the DNA. From this culture they took a sample, transferring it to a medium containing only light N, ensuring that the next round of DNA replication would have to make use of light N. If, as Crick and I had predicted, DNA replication involves unzipping the double helix and copying each strand, the resultant two "daughter" DNA molecules in the experiment would be hybrids, each consisting of one heavy N strand, the template strand derived from the "parent" molecule, and one light N strand, the one newly fabricated from the new medium. Meselson and Stahl's centrifugation procedure bore out these expectations precisely. They found three discrete bands in their centrifuge tubes, with the heavy-then-light sample halfway between the heavy-heavy and light-light samples. DNA replication works just as our model supposed it would.

\smallskip
\noindent\textbf{中文。} 他们使用一种离心技术，可以按照重量的细微差异对分子排序；离心旋转后，较重分子最终会比轻分子更靠近试管底部。氮原子（N）是DNA组成部分，而氮存在两种不同形式，一轻一重，因此梅塞尔森和斯塔尔能够标记DNA片段，从而追踪细菌中DNA复制过程。最初，所有细菌都在含重氮的培养基中培养，因此重氮被纳入DNA两条链。随后他们从该培养物中取样，转移到只含轻氮的培养基中，确保下一轮DNA复制必须使用轻氮。如果如克里克和我所预言的那样，DNA复制涉及拉开双螺旋并复制每条链，那么实验中产生的两个“子代”DNA分子将是杂合体，每个都由一条重氮链，即来自“亲本”分子的模板链，和一条轻氮链，即用新培养基新合成的链，组成。梅塞尔森和斯塔尔的离心程序精确证实了这些预期。他们在离心管中发现三条离散条带，其中先重后轻的样品位于重-重样品与轻-轻样品之间。DNA复制正如我们的模型所设想的那样运作。
\end{quote}

\begin{figure}[htbp]
\centering
\ocrimage[width=0.66\linewidth]{images/b247dd189a7c3454e931f15463bd16d36aa515299235b26261281887ab7814de.jpg}
\caption{Matt Meselson beside an ultracentrifuge, the hardware at the heart of "the most beautiful experiment in biology."\\马特·梅塞尔森站在超速离心机旁；这台设备是“生物学中最美丽的实验”的核心硬件。}
\end{figure}

\begin{quote}
\noindent\textbf{English.} The biochemical nuts and bolts of DNA replication were being analyzed at around the same time in Arthur Kornberg's laboratory at Washington University in St. Louis. By developing a new, "cell-free" system for DNA synthesis, Kornberg discovered an enzyme, DNA polymerase, that links the DNA components and makes the chemical bonds of the DNA backbone. Kornberg's enzymatic synthesis of DNA was such an unanticipated and important event that he was awarded the 1959 Nobel Prize in Physiology or Medicine, less than two years after the key experiments. After his prize was announced, Kornberg was photographed holding a copy of the double helix model I had taken to Cold Spring Harbor in 1953.

\smallskip
\noindent\textbf{中文。} 大约同一时期，圣路易斯华盛顿大学阿瑟·科恩伯格实验室正在分析DNA复制的生化细节。通过开发一种新的“无细胞”DNA合成系统，科恩伯格发现了一种酶，即DNA聚合酶，它把DNA组分连接起来，并形成DNA骨架的化学键。科恩伯格的DNA酶促合成是如此出人意料而重要，以至于关键实验后不到两年，他就获得了1959年诺贝尔生理学或医学奖。获奖公布后，科恩伯格被拍到手持一份我1953年带到冷泉港的双螺旋模型。
\end{quote}

\begin{quote}
\noindent\textbf{English.} It was not until 1962 that Francis Crick, Maurice Wilkins, and I were to receive our own Nobel Prize in Physiology or Medicine. Four years earlier, Rosalind Franklin had died of ovarian cancer at the tragically young age of thirty-seven. Before then Crick had become a close colleague and a real friend of Franklin's. Following the two operations that would fail to stem the advance of her cancer, Franklin convalesced with Crick and his wife, Odile, in Cambridge.

\smallskip
\noindent\textbf{中文。} 直到1962年，弗朗西斯·克里克、莫里斯·威尔金斯和我才获得我们自己的诺贝尔生理学或医学奖。四年前，罗莎琳德·富兰克林因卵巢癌去世，年仅三十七岁，令人痛惜。在那之前，克里克已成为富兰克林亲近的同事和真正的朋友。在两次未能阻止癌症进展的手术之后，富兰克林在剑桥与克里克及其妻子奥迪尔一起休养。
\end{quote}

\begin{figure}[htbp]
\centering
\ocrimage[width=0.62\linewidth]{images/64e65b4d42b8972cb02eb887296fc9edba05b01aeb760bb5593aa27a5d31f6b8.jpg}
\caption{Arthur Kornberg at the time of winning his Nobel Prize.\\阿瑟·科恩伯格获诺贝尔奖时。}
\end{figure}

\begin{quote}
\noindent\textbf{English.} It was and remains a long-standing rule of the Nobel Committee never to split a single prize more than three ways. Had Franklin lived, the problem would have arisen whether to bestow the award upon her or Maurice Wilkins. The Swedes might have resolved the dilemma by awarding them both the Nobel Prize in Chemistry that year. Instead, it went to Max Perutz and John Kendrew, who had elucidated the three-dimensional structures of hemoglobin and myoglobin respectively.

\smallskip
\noindent\textbf{中文。} 诺贝尔委员会长期以来一直有一条规则，且至今如此：单项奖不能分给超过三个人。如果富兰克林还活着，就会出现该把奖授予她还是授予莫里斯·威尔金斯的问题。瑞典人或许可以通过把那年的诺贝尔化学奖同时授予他们二人来解决困境。但事实上，化学奖授予了马克斯·佩鲁茨和约翰·肯德鲁，他们分别阐明了血红蛋白和肌红蛋白的三维结构。
\end{quote}

\begin{quote}
\noindent\textbf{English.} The discovery of the double helix sounded the death knell for vitalism. Serious scientists, even those religiously inclined, realized that a complete understanding of life would not require the revelation of new laws of nature. Life was just a matter of physics and chemistry, albeit exquisitely organized physics and chemistry. The immediate task ahead would be to figure out how the DNA-encoded script of life went about its work. How does the molecular machinery of cells read the messages of DNA molecules? As the next chapter will reveal, the unexpected complexity of the reading mechanism led to profound insights into how life first came about.

\smallskip
\noindent\textbf{中文。} 双螺旋的发现为活力论敲响了丧钟。严肃科学家，即使是有宗教倾向的科学家，也认识到，对生命的完整理解并不需要揭示新的自然法则。生命只是物理和化学的问题，尽管是组织得极其精巧的物理和化学。眼前的任务是弄清由DNA编码的生命脚本如何发挥作用。细胞的分子机器怎样读取DNA分子的信息？正如下章将揭示的，读取机制出人意料的复杂性，带来了关于生命最初如何出现的深刻洞见。
\end{quote}

\begin{figure}[htbp]
\centering
\ocrimage[width=0.30\linewidth]{images/dc1e2620c73dd956aa24b0dee5de0c025523e58fa15db48333e2b1f5e14dd38a.jpg}
\hfill
\ocrimage[width=0.30\linewidth]{images/28ba4827aae4f32c1a21d821087caa028bb2cd3bbbe88b6ca80cb819202d9c1f.jpg}
\hfill
\ocrimage[width=0.30\linewidth]{images/d48c5f92ded01c4d69e1782a28c8bfeeb9b00a2845d59755a7d38fec885397bf.jpg}
\caption{Uncaptioned terminal image references retained from the Watson/Berry source boundary.\\保留Watson/Berry源文本边界处未配图注的篇末图像引用。}
\end{figure}
